\section{Finite endpoint configurations}
\label{rw:finite-local}

A finite endpoint can arise when a lower scalar factor loses positivity,
when the upper arc approaches a saddle, or when these events occur
together. We first determine which combinations can be limits of actual
cycles. The exact affine upper return then gives two action columns on
the regular interval; the remaining configurations require comparisons
in the moving lower deficits and the actual upper saddle ratio.

All comparisons retain the original central height. We fix the sections
and coefficient boxes before the scale cutoff. An actual return must
traverse each whole slow path and regular connection and have positive
complete central outputs. The larger signed domains are used for
coefficient division. The local counts proved here are inputs to the
interval argument in Section~12: the three-saddle chart initially allows
four closing zeros, and its derivative estimate must be combined with
an entrance sign to obtain the sharp bound.

Table~\ref{rw:atlas:reading-map} gives the inputs and passing conditions
for the endpoint comparisons in Sections~7--10. In each row $K$
compares derivatives of the original central height. The scalar tests
hold along the whole entry--exit path; regular connections and positive
complete central outputs are also required.

\begin{center}\small
\setlength{\tabcolsep}{4pt}
\begin{longtable}{@{}p{.17\linewidth}p{.20\linewidth}p{.28\linewidth}p{.27\linewidth}@{}}
\caption{Inputs and passing tests for the endpoint comparisons.}
\label{rw:atlas:reading-map}\\
\toprule Configuration & Input outward & Passing test & Role in the final count\\\midrule
\endfirsthead
\toprule Configuration & Input outward & Passing test & Role in the final count (continued)\\\midrule
\endhead
\bottomrule\endfoot
Finite regular & Actual outgoing label $p\uparrow$ & Both lower scalar
factors positive; compact regular upper arc & At most one simple crossing of
$K$; its departure sign enters the terminal piece.\\\addlinespace
Two lower roots & Strong deficit $t\downarrow$ & Both lower deficits
positive & At most one zero of the cleared derivative; ordinary or outward
Rolle in Section~12.\\\addlinespace
Finite upper saddle & Saddle deficit $t\downarrow$ & Regular lower legs,
or each retained lower-loss passage & Actual ratio $r$; a regular-lower
entrance uses Proposition~\ref{rw:anchors:regular-saddle}.\\\addlinespace
Three finite saddles & $t\downarrow$, with the two transported lower
deficits & All three passages positive & The preliminary four-cycle
bound is sharpened by the seam signs and endpoint parity in
Proposition~\ref{rw:anchors:three-saddle}.\\\addlinespace
One or two infinite ends & Finite deficit, weak $u\downarrow$, or a
reciprocal input, as specified in Section~8 & Strict, repeated or weak
upper test and the inherited lower tests, including cubic horizontal
factors & A strict
derivative, a constant-sign connection, or the weak entrance of
Proposition~\ref{rw:anchors:primary-weak}.\\\addlinespace
Joint repeated center & Contact $p\uparrow$; separated $s\downarrow$
with $s\to1/p$ & Contact signed-gap passages; both gaps positive on the
separated part & One critical-point estimate on the contact--separated
union in Proposition~\ref{rw:joint:one-critical}.\\\addlinespace
Secondary end & Actual tag $P\uparrow$; saddle $t\downarrow$ or weak
$u\downarrow$ & Actual upper first return and inherited lower passage
on the normalized $(C,M)$ face & The regular chain feeds the saddle or
weak entrances in Section~\ref{rw:anchors:secondary}.\\
\end{longtable}
\end{center}

The regular finite generators are $(\tau_1,\sigma_1)$, whereas a
regular-lower saddle uses $(\bar a,r-1)$. The three-saddle center
near ratio one uses $(\xi_2,(r-1)\xi_1)$; away from one it uses
$(u_0,z_0)$. The joint block uses $(\tau,\zeta)$;
the secondary regular comparison uses $(\bar a,C)$ with
$C=\bar d/\lambda$. Each pair belongs to its own coefficient chart.
The local critical numerator $M$ is a positive multiple of an input
derivative of $K$.
The secondary coefficient $M$ in the normalized face $(C,M)$ has a
different meaning.

\subsection{The upper return and the necessary lower geometry}
\label{rw:finite:geometry}

The first task is to locate the possible endpoints of returning cycle
data. The zero-scale upper return determines the regular input range.
Whole-path positivity of the two lower factors and an action inequality
then exclude some combinations of losses. Use the fixed-field direction
$\ell\ge0$ chosen in Section~\ref{rw:geometry}, and put
\[
 F_\pm(X)=1\pm\bar aX+\bar BX^2,\qquad
 A_\pm(P)=\int_0^P\frac{X}{F_\pm(X)}\,\dd X
\]
on intervals where the denominator is positive throughout. Approaching
a root is allowed only from this passing side; the action is never
continued through it.

At $\epsilon=\nu=0$ the upper field, divided by its positive factor
$Y$, is
\begin{equation}\label{rw:finite:principal-upper}
 X'=\bar mY-1,\qquad Y'=X+\bar cY.
\end{equation}
Let $S''-\bar cS'-\bar mS=0$, $S(0)=0$, $S'(0)=1$, and
$I(T)=\int_0^T S(t)\,\dd t$. An arc beginning at $(p,0)$ has
$Y(t)=pS(t)-I(t)$ and $X=Y'-\bar cY$.
For $\mathcal D=\bar c^2+4\bar m\ge0$ write
$\gamma_\pm=(\bar c\pm\sqrt{\mathcal D})/2$, and define
\[
 p_c=\begin{cases}1/\gamma_+,&\mathcal D\ge0,\ \gamma_+>0,\\
 \infty,&\text{otherwise},\end{cases}
 \qquad q_c=p_c(\bar m,-\bar c).
\]

\begin{lemma}[The regular upper arc]\label{rw:finite:upper-return}
The first return of \eqref{rw:finite:principal-upper} through $Y>0$
exists exactly for $0<p<p_c$. Its exit $(-q,0)$ defines an increasing
analytic bijection $q=\mathcal U_0(p)$ from $(0,p_c)$ onto $(0,q_c)$.
If $T$ is the flight time, then
\begin{equation}\label{rw:finite:upper-derivative}
 \mathcal U_0'(p)=\frac pq e^{\bar cT},\qquad
 \mathcal U_0^{-1}(q;\bar m,\bar c)
       =\mathcal U_0(q;\bar m,-\bar c).
\end{equation}
On every compact regular coefficient and input set the positive-scale
return exists on a larger common tube and satisfies
\begin{equation}\label{rw:finite:regular-upper-error}
 \mathcal U_{\epsilon,\nu}(p)
 =\mathcal U_0(p)+O_N(\epsilon(1+|\log\epsilon|))+\nu O_N(1).
\end{equation}
Here and below $O_N$ includes selected finite ordinary input and
coefficient derivatives and Euler scale derivatives, with a finite
total order chosen before the domains. Common smaller image intervals
have bounded inverse jets in the same class.
\end{lemma}
\begin{proof}
On the first positive interval of $S$ put $W=S^2-IS'$.
Since $W'=\bar cW+S$,
\[
 W=e^{\bar cT}\int_0^T e^{-\bar ct}S(t)\,\dd t>0,
 \qquad (I/S)'=W/S^2>0.
\]
Thus the flight-time equation $p=I(T)/S(T)$ has at most one
solution. At that solution $Y'(T)=-W/S<0$, so the exit is
$(-q,0)$ with $q=W/S>0$.
For $\mathcal D<0$, $S=e^{\bar cT/2}\sin(\sqrt{-\mathcal D}T/2)
/(\sqrt{-\mathcal D}/2)$; the quotient $I/S$ tends to infinity
at the first zero of $S$. For $\mathcal D>0$ use
$S=(e^{\gamma_+T}-e^{\gamma_-T})/(\gamma_+-\gamma_-)$.
Its quotient tends to $1/\gamma_+$ if $\gamma_+>0$ and to infinity
otherwise. At $\mathcal D=0$, $S=Te^{\bar cT/2}$ gives the same
rule, including $\bar c=0$. These are exactly the stated ranges.
The determinant formula for transverse sections has normal factors
$p,q$ and divergence $\bar c$, proving the derivative identity.
Reflection gives its inverse and the range of $q$.

To estimate the positive-scale return on a common tube, split each
compact reference arc into two layers $0\le Y\le d_0$ and a middle
arc with $Y\ge d_0$.
On a layer $G=X+\bar cY$ has a fixed nonzero margin. Set
$g=\bar mY-1$, $F=\bar BX^2+\bar aX$,
$Q=\epsilon/(Y+\epsilon)$ and $r_0=Y^2/(Y+\epsilon)$.
The exact equation there is
\[
 X_Y=\frac{(1-Q)g+QF}{G+\nu r_0}.
\]
The perturbation at $\nu=0$ is $Q(F-g)/G$, whose integral is
$O(\epsilon\log(1+d_0/\epsilon))$. In the middle, use the
positive time change by $Y+\epsilon$; its field is
$(g+Q(F-g),G+\nu r_0)$. Ordinary flow dependence and the transverse
exit equation give the value estimate on larger compact tubes.
For the selected jets use the bounded-layer variation argument
\eqref{rw:app:bounded-variation}, with $E_\epsilon=\epsilon\partial_\epsilon$
and the integrable carriers $E_\epsilon Q=Q(1-Q)$,
$E_\epsilon r_0=-Qr_0$; further Euler
rows retain these carriers times bounded polynomials in $Q$.
Ordinary $\nu$ integration gives the last term of
\eqref{rw:finite:regular-upper-error}. The positive compact margin
in \eqref{rw:finite:upper-derivative} then gives the inverse jets on
fixed smaller common images.
\end{proof}

The lower maps must be expressed in the original central height to
compare the two complete paths. Their absolute size is exponentially
small, so we estimate a logarithmic amplitude after removing the action.
The common parameter domains below also contain the center contractions
of Section~5.

\begin{lemma}[Regular lower amplitudes]\label{rw:finite:regular-lower}
Fix a compact set of positive finite inputs with $F_+>0$ on $[0,p]$.
For $p$ bounded away from zero, the signed outgoing backward map
from $(p,1)$ to $X=0$ is
\begin{equation}\label{rw:finite:lower-height}
 L_+(p)=Z_+(p)-\ell H_+(p,\ell),\qquad
 Z_+(p)=a_+(p)e^{-A_+(p)/\epsilon},
\end{equation}
on a common enlarged input and coefficient domain. The logarithm of
$a_+>0$ has bounded selected jets, while
$H_+=\sqrt{\pi/2}+O_N(\sqrt\epsilon+|\ell|)$ and its required
ordinary $\ell$ derivatives are bounded. Reflection gives
$L_-(q)=Z_-(q)+\ell H_-(q,\ell)$ with the same properties.
Positive endpoints of either signed map imply a passage with $z>0$.

With the fixed weak references of Section~4 one can instead write
\begin{equation}\label{rw:finite:relative-lower}
 L_+(p)=s_++U(p),\quad L_-(q)=s_-+V_0(q),\qquad
 U=\alpha_Ue^{-A_+(p)/\epsilon},\quad
 V_0=\alpha_Ve^{-A_-(q)/\epsilon}.
\end{equation}
The positive amplitudes have bounded logarithmic jets on a small
ordinary $\nu$ box, and $s_--s_+=\rho\nu H$, $H>h_*>0$,
is independent of input. In particular $U',V_0'<0$ and their
logarithmic derivative rates have the usual action limits on every
smaller compact regular input interval.
\end{lemma}
\begin{proof}
Choose a small fixed $d_0>0$. On the regular layer $d_0\le z\le1$,
use $z$ as independent variable. Its denominator has a nonzero margin,
so $P=p+\epsilon R$ has bounded selected jets.
On the core use \eqref{rw:relative:core-field}, with
$\nu=\ell/\rho$ and $F_\epsilon=F_++\epsilon\bar m$.
Its linear part is
$(X/(\epsilon F_\epsilon)-\bar c/F_\epsilon)z
 +\ell/(\rho F_\epsilon)$.
The backward kernel $K(X,u)$ on $0\le X\le u\le P$ satisfies
\[
 |K(X,u)|\le Ce^{-c(u^2-X^2)/\epsilon},\quad
 \int_X^P|K|\frac u\epsilon\,\dd u\le C,
 \quad \int_X^P|K|\frac{\dd u}{\rho}\le C.
\]
The remainder has size
$C((X/\epsilon+1)z^2+|\ell z|/\rho)$.
Choose $d_0$ and then a fixed ordinary $\ell$ disk so that the
backward integral is a contraction. This constructs the signed map
as well. On $z=0$ the graph derivative has the sign of $\ell$, so
a trajectory with positive endpoints cannot leave and reenter $z>0$.
At $\ell=0$ the line itself is invariant.

At zero forcing,
$z(X)\le Cd_0e^{-c(P^2-X^2)/\epsilon}$, so
$\int_0^P(X/\epsilon+1)z(X)\,\dd X\le Cd_0$.
Removing the action leaves a bounded logarithm: the endpoint shift
$(A_+(P)-A_+(p))/\epsilon$ is bounded since $P-p=\epsilon R$,
the linear correction uses $F_\epsilon-F_+=O(\epsilon)$, and the
nonlinear correction is the preceding integral.
The normalized Volterra argument of
Proposition~\ref{rw:transitions:central-core} supplies its selected
jets after $X=Py$: $P$ stays in a compact positive interval, its
regular endpoint jets are bounded, and the kernel satisfies the same
weighted bound \eqref{rw:app:normalized-kernel-bound}.

Differentiate the core integral with respect to $\ell$ at zero
forcing. Up to an exponentially small layer endpoint term, its
negative outgoing value is
\[
 \int_0^P
 \exp\left[-\int_0^X \mathfrak f_z(u,z_0(u))\,\dd u\right]
 \frac{(1-z_0(X))^2}{\rho D(X,z_0(X))}\,\dd X.
\]
Near zero put $X=\rho v$. The integrand tends to $e^{-v^2/2}$;
its error is $O(\rho)$ times a fixed polynomial and a Gaussian.
Away from a fixed small central interval the positive linear rate
makes the error exponentially small. The same kernels bound higher
ordinary $\ell$ variations. Hadamard integration therefore proves
the formula for $H_+$ and its jets. Reflection proves the incoming
formula.

For \eqref{rw:finite:relative-lower}, restrict the one chosen weak
reference $z_b$ at the variable endpoint $P$. Its continuation to
$X=0$ is the same reference for all inputs. On the weighted Gaussian
core of Proposition~\ref{rw:transitions:central-core}, the difference
$w=z-z_b$ satisfies $w_X=a(X)w+b(X,w)w^2$ with
\[
 a-X/(\epsilon F_+)=O(1),\quad
 \partial_\nu a=O(1),\quad |b|\le C(X/\epsilon+1).
\]
The weighted reference estimate gives
$(X/\epsilon)|\partial_\nu z_b|=O(1)$, controlling the potentially
large coefficient term. It is this weighted bound that is needed here,
rather than bounds on the unweighted high $X$ derivatives.
The positive bounded endpoint deviation and the same normalized
Volterra estimate give $\alpha_U$; reflection gives $\alpha_V$.
The Gaussian variation above, now with $\ell=\rho\nu$, gives
opposite reference signs and gap
$\sqrt\epsilon\,\nu(\sqrt{2\pi}+O(\sqrt\epsilon))$.
Finally $-\epsilon U'/U=A_+'-\epsilon(\log\alpha_U)'$ and its
reflected identity give the positive derivative margins and logarithmic
rates, using relative rather than absolute errors.
\end{proof}

\begin{lemma}[Entry bounds and action ordering]
\label{rw:finite:action-order}
There is a fixed $\chi>0$ such that every parameter-dominant cycle
in the chosen direction has $p>\chi/2$ and $q>\chi$ at small scale.
If its incoming data stay in a compact set where $F_->0$ on $[0,q]$,
then its outgoing data stay in a compact regular lower set and
\begin{equation}\label{rw:finite:action-inequality}
 A_+(p)\le A_-(q)+C\epsilon.
\end{equation}
Consequently outgoing loss with a regular incoming lower limit is
impossible.
\end{lemma}
\begin{proof}
Take $\chi=\eta/4$, with the fixed energy threshold of Section~2,
and decrease $\eta\le1/16$ if necessary. The crossing argument in
that section gives one intersection with each half of $y=0$.
For $0\le X\le\chi$, $0\le z\le1/8$, the lower denominator
is greater than $3/4$. The cycle has $0\le\ell<3s$ and
$s<e^{-\eta^2/\epsilon}$, so variation of constants gives
\[
 z_X\le(4X/(3\epsilon)+4/3)z+4s/\rho,
 \qquad z(X)\le6s e^{2X^2/(3\epsilon)+4X/3}
                  <e^{-\eta^2/(2\epsilon)}.
\]
Thus the orbit reaches $X=\chi$. A later decrease to $\chi/2$
cannot occur while $z\le1/2$, where its horizontal velocity is
positive. In $1/2<z<1$, $z_\xi\ge\chi z/4$ as long as
$X\ge\chi/2$, and the possible decrease in $X$ is at most
$4\epsilon\log2/\chi<\chi/2$. Hence $p>\chi/2$.
For the incoming side, the reference arc with $q=\chi$ has
$F_->3/4$ and $A_-(\chi)\le\eta^2/24$. Lemma~\ref{rw:finite:regular-lower}
gives $s_{\rm ref}\ge c e^{-\eta^2/(24\epsilon)}>s$.
This arc, the intervening central segment and the segment on $z=1$
bound a flow region: the top is an entry, the central edge an exit,
and uniqueness prevents crossing the reference arc. Every $q\le\chi$
would consequently give height at least $s_{\rm ref}$, a contradiction.

Now fix incoming compact regular data, so
$s\ge ce^{-A_{\max}/\epsilon}$ for a fixed $A_{\max}$.
On a fixed outgoing strip $0<X\le R$, a turn $D=0$ is a strict
maximum of $X$. Indeed, with $u=1-z$ its equation gives
$u(1+\epsilon\bar m u)=-X(\bar a+\bar BX)$ and
$u\le2(1+R)X$. Thus $z_\xi>0$ and $D_\xi=D_z z_\xi<0$
at the turn, for small scale. Before a first rightward crossing of
$X=R$ one has $D>0$. Also $F_\epsilon-D>0$, and the forcing
term is nonnegative. On $[0,\chi]$ use
$z_X/z\ge X/(\epsilon F_\epsilon)-4/3$, and beyond $\chi$ use
$z_X/z\ge(1-\epsilon/\chi)X/(\epsilon F_\epsilon)$.
Since $z(R)\le1$, integration gives
\begin{equation}\label{rw:finite:action-budget}
 (1-\epsilon/\chi)\int_0^R\frac{X}{F_\epsilon(X)}\,\dd X
       \le-\epsilon\log s+4\chi\epsilon/3.
\end{equation}
The incoming regular arc bounds the right side by $A_{\max}+o(1)$.
Since $F_\epsilon(X)\le2+X+X^2$, the integral on the left is at
least $\frac12\log(R/2)$ for $R\ge2$. A fixed sufficiently large $R$
therefore bounds $p$.
If $F_+$ then approached zero at or before a limiting endpoint,
apply the same inequality for every $R$ before its first zero and
let $R$ approach that zero. The integral diverges for both simple
and double zeros, contradicting the bound. This argument also
excludes nearby positive polynomials with minima tending to zero.
Thus both lower data sets are compact and regular. At a closing zero,
$L_+(p)=L_-(q)$ and $\ell\ge0$ give $Z_+(p)\ge Z_-(q)$.
Taking logarithms with their bounded amplitudes proves
\eqref{rw:finite:action-inequality}.
\end{proof}

\begin{lemma}[Lower cones and simultaneous finite loss]
\label{rw:finite:lower-cones}
At a limit $(\mathbf b_*,p_*,q_*)$ of cycle data, $F_+$ and $F_-$
are nonnegative up to their respective endpoints. In particular
\[
 p_* =\infty\Longrightarrow\bar B\ge0,\ \bar a\ge-2\sqrt{\bar B},
 \qquad
 q_* =\infty\Longrightarrow\bar B\ge0,\ \bar a\le2\sqrt{\bar B}.
\]
If both finite lower limits lose positivity, then $\bar B=-\beta<0$
and their endpoints are the simple roots
\[
 p_*=(\bar a+\Delta)/(2\beta),\qquad
 q_*=(-\bar a+\Delta)/(2\beta),\qquad
 \Delta=\sqrt{\bar a^2+4\beta},\qquad p_*q_*=1/\beta.
\]
If the upper limit is regular, these roots satisfy
$\mathcal U_0(p_*)=q_*$.
\end{lemma}
\begin{proof}
At an outgoing line $X=R>0$, if $F_+(R)+\epsilon\bar m\le0$,
the horizontal velocity is negative for $0<z<1$. It prevents the
first rightward crossing needed to reach $p>R$. Reflection gives
the corresponding incoming obstruction. Passing to the limit proves
nonnegativity and the two quadratic cone conditions.
For $\bar B\ge0$, a positive root of $F_+$ requires $\bar a<0$,
whereas one of $F_-$ requires $\bar a>0$. Hence both losses force
$\bar B<0$ and the displayed roots are unique and simple. The
same barrier for the nearby exact quadratic roots prevents a cycle
endpoint from lying beyond either root. Since each limiting root
was assumed at or before its endpoint, they coincide. The upper
identity follows from the regular return convergence.
\end{proof}

A strict real upper threshold also excludes returns beyond it,
uniformly for unbounded inputs. We will use this fact at infinity.
If $\mathcal D>0$, $\gamma_+>0$ on a compact coefficient block, then for every fixed $\zeta>0$ no input
$p\ge p_c+\zeta$ returns, at a sufficiently small common scale.
To prove this without bounding $p$, choose $r>0$ between the two
eigenvalues and close to $\gamma_+$, and choose $b_0>1/r$ with
$b_0<p_c+2\zeta/3$. Then
$A=\bar m+r\bar c-r^2>0$ and $rb_0-1>0$ have fixed margins.
The limiting reference arc from $p_c+\zeta$ reaches a small
$Y=d_0$ in $X+rY-b_0>0$. On the sloping boundary of the unbounded
region $Y\ge d_0$, $X+rY-b_0\ge0$, its exact outward derivative
in the time divided by $Y+\epsilon$ is
\[
 AY+rb_0-1+Q\{F(b_0-rY)-g(Y)\}
                         +r\nu Y^2/(Y+\epsilon)>0.
\]
The error is bounded by $C_{d_0}\epsilon(1+Y)+C|\nu|Y$;
the bottom boundary points inward as well. This is a forward
invariant region. Larger inputs either reach it, by uniqueness of
the layer graphs, or escape affinely; neither returns. Reflection
gives the exit exclusion. This proof requires a strict discriminant:
it gives no exclusion at a repeated root.
\label{rw:finite:strict-barrier}

Thus two infinite endpoints require $\bar c^2+4\bar m\le0$,
because otherwise at least one reflected threshold is finite.
The remaining upper limits have $\bar m<0$, or
$\bar c=\bar m=0$. In the latter case the lower cones and primary
normalization imply $\bar B>0$. If only $p_*$ is infinite and
$q_*$ finite, the incoming lower polynomial must lose positivity,
by Lemma~\ref{rw:finite:action-order}. These restrictions classify
limits of actual cycle data. The passage constructions still use
adjacent signed coefficient domains for their common comparisons.

\subsection{Two columns on the regular interval}

On a compact regular interval, the two center generators give two
positive action columns with a strictly monotone quotient. This controls
both the number of critical zeros and their crossing direction. We first
state the elementary perturbation argument used to pass from the action
to the actual return.

\begin{lemma}[Two columns and their crossing direction]
\label{rw:atlas:two-columns}
On one interval let $A,B$ have uniformly bounded $C^1$ norms, with
$|B|>b_*>0$ and $(A/B)'$ of constant sign and magnitude greater
than $r_*>0$. Put $Q=aA+bB+e$, where $a,b$ are constant,
$I=|a|+|b|>0$ and $\|e\|_{C^1}\le\eta I$.
For a sufficiently small common $\eta$, every zero is simple,
$|a|\ge cI$ there, and its derivative has sign
$aB(A/B)'$. Thus $Q$ has at most one zero, including multiplicity.
If moreover $A,B,A',B'$ and $(A/B)'$ have positive margins, then
$Q'$ has sign $a$ wherever $Q=0$ or $Q$ has sign $a$.
When $a=0$, both $Q$ and $Q'$ have sign $b$.
\end{lemma}
\begin{proof}
At a zero, $a(A/B)+b=-e/B$, so the bounded ratio and nonvanishing
$B$ give $|a|\ge cI$. Differentiating at that zero gives
\[
 Q'=aB(A/B)'+e'-B'e/B.
\]
The last two terms are $O(\eta I)$, proving the strict direction.
Two crossings with this same direction are impossible.
For the additional assertion, if $a,b$ have the same sign, use the
positive column derivatives. If their signs are opposite, at a point
where $Q=0$ or $Q$ has sign $a$,
$\operatorname{sign}(a)[a(A/B)+b]\ge-|e|/B$, again forcing
$|a|\ge cI$. In
$Q'=B'[a(A/B)+b]+aB(A/B)'+e'$ the second term has sign $a$ and dominates the possible
$O(\eta I)$ contribution of the opposite sign. The case $a=0$ follows from
the margins of $B,B'$.
\end{proof}

We keep the original central height as output. Write
\[
 \tau=\bar a+\bar c,\qquad
 \varsigma=\bar B+\bar m,\qquad \sigma=\bar c\varsigma.
\]
Intersect the upper return range $0<p<p_c$ with $F_+>0$ on
$[0,p]$ and $F_->0$ on $[0,\mathcal U_0(p)]$. Since
$\mathcal U_0$ increases, this regular limiting domain is one interval
with left endpoint zero. On it set
$\mathcal A=A_-(\mathcal U_0(p))-A_+(p)$.
A polynomial quasi-invariant of the affine upper field expresses
$\mathcal A$ in the two center generators. Its columns are positive
integrals, and two weighted-average arguments prove that their quotient
increases. The same calculation will determine the direction in which
the actual critical function crosses zero.

\begin{lemma}[The regular action columns]\label{rw:regular-columns}
For the upper orbit of flight time $T$ in
Lemma~\ref{rw:finite:upper-return}, define
\[
 J_1(T)=\int_0^T e^{-\bar ct}(1-\bar mY(t))\,\dd t,
 \qquad J_2(T)=\int_0^T e^{-\bar ct}Y(t)\,\dd t.
\]
Then $J_1=p+e^{-\bar cT}q>0$, $J_2>0$, and
\begin{equation}\label{rw:finite:quasi-invariant}
 F_-(q)-e^{\bar cT}F_+(p)
       =-e^{\bar cT}(\tau J_1+\sigma J_2).
\end{equation}
On the regular interval,
\begin{equation}\label{rw:finite:action-columns}
 \mathcal A=\tau A+\sigma C,\qquad
 A=\int_0^T\omega J_1\,\dd t,
 \quad C=\int_0^T\omega J_2\,\dd t,
 \quad \omega=\frac{p p_Te^{\bar cT}}{F_+(p)F_-(q)}>0.
\end{equation}
Both columns and their $p$ derivatives are positive, and
$(C/A)'>0$. Every nonzero constant combination of $A,C$ has at
most one zero, which is simple. At the small endpoint, uniformly
on bounded normalized coefficient sets,
\begin{equation}\label{rw:finite:small-columns}
 A=\tfrac23p^3(1+O(p)),\qquad
 C=\tfrac2{15}p^5(1+O(p)),\qquad C/A=\tfrac15p^2(1+O(p)).
\end{equation}
\end{lemma}
\begin{proof}
Introduce the polynomial
\[
 V=1+\bar aX+\bar BX^2+(2\bar B+\bar c\bar a)Y
                  +\bar c\bar BXY-\bar B\bar mY^2.
\]
Substitution of \eqref{rw:finite:principal-upper}, grouping its
constant and $Y$ terms, gives
$V'-\bar cV=-\tau(1-\bar mY)-\sigma Y$.
Its endpoint values are $F_+(p)$ and $F_-(q)$, so integration
proves \eqref{rw:finite:quasi-invariant}.
The upper arc has $Y>0$ in its interior, giving $J_2>0$.
Using $1-\bar mY=-X'$ and
$e^{-\bar ct}X=(e^{-\bar ct}Y)'$, integration by parts gives
$J_1=p+e^{-\bar cT}q$.

For monotonicity, put $u_1(T)=\int_0^T e^{-\bar ct}\,\dd t$.
Then $J_1=u_1-\bar mJ_2$.
From $Y=p(T)S-I$ and the formulas for $W,p_T,q$ above,
\[
 J_2'=p_T\int_0^T e^{-\bar ct}S(t)\,\dd t
       =e^{-\bar cT}W^2/S^2=e^{-\bar cT}q^2.
\]
Thus $J_2'/u_1'=q^2$ is strictly increasing. Since both integrals
vanish at zero, $J_2/u_1$ is its weighted average and is strictly
smaller than its terminal value. Hence
\[
 (J_2/u_1)'=\frac{u_1'}{u_1}(q^2-J_2/u_1)>0,
 \qquad
 (J_2/J_1)'=
 \frac{(J_2/u_1)'}{(1-\bar mJ_2/u_1)^2}>0.
\]
Here $1-\bar mJ_2/u_1=J_1/u_1>0$.

Differentiate $A_-(q)-A_+(p)$ and use
$q_T=(p/q)e^{\bar cT}p_T$. The result is
\[
 \mathcal A_T=
 \frac{p p_T}{F_+(p)}
       \left(\frac{e^{\bar cT}F_+(p)}{F_-(q)}-1\right)
             =\omega(\tau J_1+\sigma J_2).
\]
Both actions vanish at zero, so this proves
\eqref{rw:finite:action-columns}. The ratio $C/A$ is the weighted
average of the strictly increasing $J_2/J_1$ with positive weight
$\omega J_1$. The same averaging argument gives $(C/A)_T>0$,
hence also $(C/A)_p>0$. It proves the zero assertion directly;
positivity of the derivatives follows from the positive integrands.
Finally scale $X=p\xi$, $Y=p^2\eta$, $t=p\theta$ in the
affine field. It tends uniformly to
$\xi'=-1$, $\eta'=\xi$, whose arc is
$\eta=(1-\xi^2)/2$. Thus
$T=2p(1+O(p))$, $J_1=2p(1+O(p))$,
$J_2=\frac23p^3(1+O(p))$ and $\omega=p/2(1+O(p))$.
Integration with $\dd T=2(1+O(p))\dd p$ gives
\eqref{rw:finite:small-columns}.
\end{proof}

\begin{proposition}[The regular critical function]
\label{rw:finite:regular-count}
Fix compact regular upper and lower data, with the input bounded
away from zero. On common slightly larger tubes the closing function
is exactly
\[
 F(p)=U(p)-V(p)-\rho\nu H,
 \qquad V=V_0\circ\mathcal U_{\epsilon,\nu},
 \qquad K=\epsilon\log(V'/U').
\]
Here $U',V'<0$, and $H>h_*>0$ is input independent.
There are uniform choices of tubes, $\theta_0>0$ and a scale cutoff
such that, on any one interval contained in the tubes, if
\begin{equation}\label{rw:finite:regular-forcing}
 |\nu|\le\theta_0(|\tau|+|\sigma|),
\end{equation}
then either the field is a center or $K$ has at most one zero,
which is simple. Consequently $F$ has at most two isolated zeros,
counting multiplicity, on the positive central passing interval.
The coefficient condition holds with a factor tending to zero in
place of $\theta_0$ whenever the fibre contains a closing zero
in one of these compact regular tubes.
\end{proposition}
\begin{proof}
Use the references of Lemma~\ref{rw:finite:regular-lower}, fixed
before the input is varied. Put $q=\mathcal U_{\epsilon,\nu}(p)$
and
\[
 r_U=A_+'-\epsilon(\log\alpha_U)',\qquad
 r_V=A_-'-\epsilon(\log\alpha_V)'.
\]
They have positive margins. Direct differentiation gives the exact
identity
\begin{equation}\label{rw:finite:regular-logarithm}
 -K=A_-(q)-A_+(p)
       +\epsilon\log\frac{\alpha_U(p)r_U(p)}
        {\alpha_V(q)r_V(q)\mathcal U_{\epsilon,\nu}'(p)}.
\end{equation}
The argument of the logarithm is positive and its logarithm has
bounded selected finite jets. Thus the right side converges, with
those jets, to $\mathcal A$ at $\nu=0$, with error
$O(\epsilon(1+|\log\epsilon|))$; changing $\nu$ adds $O(\nu)$.

The real period domains of Section~5 contain each compact regular
center tube under consideration. On the second component, the
separated coordinate at $Y=0$ is $\chi=-p$ and
$A_C(-p)=F_+(p)+O(\epsilon)$ throughout the bounded interval.
Its positive margin places the entire arc in the actual center
domain. On the reversible component the compact upper arc lies
strictly inside the central potential well. Hence in either case
the complete central maps agree. Since the reference values vanish
at $\nu=0$, $K=0$ there. The zero-scale action has the same two
identities by \eqref{rw:finite:action-columns}.

Choose a common larger regular coefficient domain near each
center tube and apply Proposition~\ref{rw:generator-division} to
the difference between \eqref{rw:finite:regular-logarithm} and
$\mathcal A$. Two additional coefficient derivatives suffice.
Away from the center set the generator norm has a positive margin,
so the error can instead be divided by that norm. On the resulting
compact cover we have
\begin{equation}\label{rw:finite:perturbed-columns}
 -K=\tau A_\epsilon+\sigma C_\epsilon+\nu E_\epsilon,
 \qquad (A_\epsilon,C_\epsilon)\longrightarrow(A,C)
                  \quad\hbox{in }C^2_p,
 \quad \|E_\epsilon\|_{C^2_p}\le C.
\end{equation}
The common center domain is essential here: it allows division of
the relative logarithmic error by the exact generators, with the
stated input-derivative bounds.

On the larger tubes both actions have a common lower bound $S_*>0$.
The amplitude and upper-map estimates bound the required absolute
coefficient derivatives by $C\epsilon^{-M}e^{-S_*/\epsilon}$,
also on the small $\nu$ box after reducing $S_*$. The same center
division preserves this weight. Lemma~\ref{rw:centers:relative-forcing},
with $I=|\tau|+|\sigma|$, therefore gives
$|\nu|=o(1)I$ at a closing zero; for $I=0$ it gives the center
identity and no isolated zero.

Choose $\theta_0$ and then the scale so that
\eqref{rw:finite:perturbed-columns} satisfies
Lemma~\ref{rw:atlas:two-columns}, with columns $(C,A)$ and
coefficients $(\sigma,\tau)$. The strict quotient margin of
Lemma~\ref{rw:regular-columns} gives at most one simple zero of $K$.
Lemma~\ref{rw:zero-calculus} gives the stated count for $F$.
Since $L_+(p)$ and $L_-(q(p))$ decrease, their positive sets are
initial subintervals. Intersecting them with the regular tube gives
one positive-central passing interval, to which the count applies.
\end{proof}

\begin{lemma}[Direction on the regular part]
\label{rw:finite:regular-direction}
Under the preceding choices, if $\sigma\ne0$, then at every point
where $K=0$ or $K$ has sign $-\sigma$, its derivative in increasing
$p$ has sign $-\sigma$. After its possible zero it therefore keeps
this sign and increases in absolute value on the rest of the regular
passing interval. If $\sigma=0$, $\tau\ne0$, both $K$ and $K'$
have sign $-\tau$.
\end{lemma}
\begin{proof}
The positive margins of $A,A',C,C'$ and $(C/A)'$ on each compact
regular tube allow the last clause of Lemma~\ref{rw:atlas:two-columns}
to be applied to $-K$ in \eqref{rw:finite:perturbed-columns}.
At a zero, $-K$ crosses into the side of sign $\sigma$, and its
derivative keeps that sign there. Thus it stays on that side. A regular
input change preserves this direction for the same original-central
function $K$.
\end{proof}

\subsection{An incoming loss and two matched lower roots}

An incoming lower loss can have unbounded transit time. To compare
its central derivative with that of a regular outgoing arc, we use
negative divergence on the outer arc and contraction through a fixed
Gaussian core. This argument also allows an incoming endpoint at
infinity. When both lower roots are lost, neither derivative provides
a regular reference; we instead compare their logarithmic rates in
the two moving deficits.

\begin{lemma}[Incoming contraction]\label{rw:finite:incoming-contraction}
Suppose the outgoing lower data are compact and regular, with
$A_+(p)\le A_{\rm out}$. Choose a fixed incoming crosscut $R>0$
with $F_->f_*>0$ on $[0,R]$ and entries $q>R+2\zeta$, where
$\zeta>0$ is fixed. On the actual common positive-central passing
set, for either sign of $\ell$, the incoming map $\mathcal I(q)$
satisfies
\begin{equation}\label{rw:finite:incoming-bound}
 0<-\mathcal I'(q)\le\frac C\epsilon e^{-A_-(R)/\epsilon}
        \left(\frac{1+R^2}{1+q^2}\right)^{1/(24\epsilon)}.
\end{equation}
If $A_-(R)>A_{\rm out}+3$ and the upper derivative is bounded,
then $K<0$ on that set. The closing function is strictly monotone
on each of its intervals. Every closing zero there is simple and
has return multiplier less than one.
\end{lemma}
\begin{proof}
In the incoming coordinates $W=-X$, the forward divergence on
$W\ge R$, $0<z<1$ is
\[
 -(1+2\epsilon\bar B)W
 +\epsilon\{\bar a+\bar c(2z-1)\}
                      -2\rho\ell(1-z)\le-W/2.
\]
The normal component satisfies
$z'\le-(R/2)z+\rho\max(\ell,0)$.
On the fixed strip $R\le W\le R+\zeta$ the horizontal speed
has magnitude at most $M\epsilon$. Any passing arc therefore
reaches $W=R$ with
\begin{equation}\label{rw:finite:height-at-crosscut}
 z_R\le e^{-R\zeta/(2M\epsilon)}
                          +2\rho\max(\ell,0)/R.
\end{equation}
An arc that has already reached $z=0$ is outside the stated domain.

The fixed core from this crosscut to the center has denominator
\[
 D_-=F_-+\epsilon\bar m-(1+2\epsilon\bar m)z
                                      +\epsilon\bar m z^2>f_*/2
\]
on a small height strip, and its graph equation is
\[
 z_W=\frac{Wz+\epsilon\bar c z(1-z)-\rho\ell(1-z)^2}
                  {\epsilon D_-}=:f(W,z,\ell).
\]
Let its signed output be $\mathcal C_R(z_R,\ell)$.
The Gaussian estimate gives
\[
 0\le\mathcal C_R(z_R,0)\le Cz_Re^{-A_-(R)/\epsilon}.
\]
For this bound, integrate
$\partial_W\log z=W/(\epsilon F_-)+O(1+(W/\epsilon+1)z)$;
the nonlinear integral is bounded by $Cz_R$.
For $\ell<0$ with positive central output, the solutions for
$\theta\ell$, $0\le\theta\le1$, lie between this positive
solution and the unforced one by scalar comparison. At fixed $z_R$,
variation of constants gives
\[
 \partial_\ell\mathcal C_R(z_R,\theta\ell)
 =\int_0^R e^{-\int_0^u f_z(v,z(v),\theta\ell)\,\dd v}
       \frac{(1-z(u))^2}{\rho D_-(u,z(u))}\,\dd u\ge c>0.
\]
To obtain the positive lower bound, restrict the integral to
$0\le u\le\rho$. Its prefactor is at least $c/\rho$, and
$|f_z|\le C(W/\epsilon+1+|\ell|/\rho)$ bounds the exponent
there. Positivity of the final output consequently implies
$|\ell|\le Cz_Re^{-A_-(R)/\epsilon}$.
For $\ell\ge0$, positivity of the outgoing regular map instead
gives $\ell\le CZ_+(p)$. Thus $\ell$ is exponentially small
on either part of the common passing set. Equation
\eqref{rw:finite:height-at-crosscut} now makes $z_R$ exponentially
small as well.

Set $\Phi(W)=\log(1+W^2)/(12\epsilon)$.
Since $W'=-\epsilon D_-$ and
$D_-\le2+W+W^2\le3(1+W^2)$, its time derivative is at least
$-W/2$, hence at least the divergence. The determinant formula
from $(q,1)$ to $W=R$ therefore bounds the transverse derivative
by
\[
 \frac{Cq}{\epsilon}
      \left(\frac{1+R^2}{1+q^2}\right)^{1/(12\epsilon)}.
\]
Half this power absorbs $q$ uniformly even for unbounded entries.
On the core, its normal variational exponent is
$A_-(R)/\epsilon+O(1)$, because
\[
 f_z=\frac{W}{\epsilon(F_-+\epsilon\bar m)}
       +\frac{\bar c}{F_-+\epsilon\bar m}
       +O((W/\epsilon+1)|z|+|\ell|/\rho),
\]
and the last two quantities have the bounds just proved.
Composition gives \eqref{rw:finite:incoming-bound}.
The outgoing relative amplitude has
$-L_+'(p)\ge(c/\epsilon)e^{-A_{\rm out}/\epsilon}$ on the
same compact regular tube, including the small forcing just bounded.
Thus the ratio of incoming and outgoing central derivatives is
less than $1/2$ when $A_-(R)>A_{\rm out}+3$ and the upper
derivative is bounded. This proves $K<0$ and the remaining claims.
The estimate uses divergence integrated along the actual passing arc
and remains valid for arbitrarily long residence near the incoming root.
\end{proof}

At a first finite simple or double incoming root, $A_-(R)\to\infty$
as $R$ approaches the root from below. If the outgoing arc and upper
return remain regular, choose $R$ so that the preceding lemma applies.
The resulting terminal part is monotone even though its incoming
transit time diverges.

For simultaneous finite lower losses,
Lemma~\ref{rw:finite:lower-cones} gives $\bar B=-\beta<0$ and
simple roots $p_l,q_l$ with $p_lq_l=1/\beta$. The next lemma
determines when a regular upper return joins these roots and supplies
the transverse parameter needed to measure their splitting.

\begin{lemma}[The regular two-root match]
\label{rw:finite:two-root-geometry}
For fixed $\beta>0$, $\bar a$ and $\bar m$, the equation
$\mathcal U_0(p_l;\bar m,\bar c)=q_l$ has a unique regular
solution in $\bar c$ exactly when $\bar m<\beta$.
At that solution the derivative with respect to $\bar c$ is
strictly positive. If $\bar a\ne0$, then
\begin{equation}\label{rw:finite:two-root-signs}
 \bar a\bar c<0,\qquad
 \operatorname{sign}(\bar a+\bar c)=\operatorname{sign}\bar c,
 \qquad |\bar c|>|\bar a|.
\end{equation}
For $\bar a=0$ the solution is $\bar c=0$.
\end{lemma}
\begin{proof}
On a returning upper arc $\bar mY-1<0$. Indeed the first
maximum of $Y$ must have $Y''=\bar mY-1<0$; the region below
that nullcline contains the complete arc. Its graph equation is
$Y_X=(X+\bar cY)/(\bar mY-1)$. Differentiation in $\bar c$,
with integration from $X=p_l$ in the decreasing direction, has
positive forcing after reversing the integration direction, because
$Y>0$ and $\bar mY-1<0$. It gives $Y_{\bar c}>0$ at the
exit and hence $\partial_{\bar c}\mathcal U_0>0$.
For $\bar m>0$, the formulas for $S,I$ show that the range at
fixed $p_l$ is $(0,1/(\bar mp_l))$: as the input approaches its
saddle threshold, the exit approaches the reflected threshold and
their product is $1/\bar m$. For $\bar m\le0$ that range is
$(0,\infty)$, by the same exponential, repeated and trigonometric
formulas, or their limits as $\bar c$ varies. Thus $q_l=1/(\beta p_l)$
belongs to the range exactly when $\bar m<\beta$.

To determine the coefficient signs at the match, put
$M=\beta-\bar m>0$ and $\tau=\bar a+\bar c$.
The endpoint identities from the solution $pS-I$ give
\[
 p_l+q_l=S(T)p_l(Mp_l-\tau)
       =e^{-\bar cT}S(T)q_l(Mq_l+\tau),
\]
and hence
\begin{equation}\label{rw:finite:matching-tau}
 \tau=M\frac{e^{\bar cT}p_l^2-q_l^2}
                      {e^{\bar cT}p_l+q_l}.
\end{equation}
The mechanical equation $Y''-\bar cY'-\bar mY=-1$ has energy
$\frac12(Y')^2+Y-\frac12\bar mY^2$ with derivative
$\bar c(Y')^2$. Thus $q_l-p_l$ has the sign of $\bar c$.
For the sign of \eqref{rw:finite:matching-tau}, write
$S(t)=e^{\bar ct/2}S_0(t)$. The positive fundamental solution
$S_0$ is hyperbolic, linear or trigonometric. For $T/2<u<T$
one has $S_0(u)>S_0(T-u)$: in the trigonometric case use
$T<\pi/\sqrt{-\bar m-\bar c^2/4}$ and the sine difference
formula. Pairing $u$ and $T-u$ then shows that
\[
 \int_0^T e^{\bar cu/2}S_0(u)\,\dd u
       -e^{\bar cT/2}\int_0^T e^{-\bar cu/2}S_0(u)\,\dd u
\]
has sign $\bar c$.
Since $q_l/p_l=e^{\bar cT}I_{-\bar c}(T)/I_{\bar c}(T)$,
this is precisely the sign of $e^{\bar cT}p_l^2-q_l^2$.
It proves the sign of $\tau$. Finally $p_l-q_l=\bar a/\beta$
gives \eqref{rw:finite:two-root-signs} and its zero case.
\end{proof}

Use the actual strong intersections of the two lower roots, their
positive regular power inputs, and the regular upper map. Common
finite-jet convergence preserves the preceding transverse derivative,
so their splitting is a coefficient coordinate $\delta$. The connector
then has the form $g(t)=(t-\delta)k(t)$, with $k,g'>c>0$ and
bounded selected jets on larger signed intervals. The deficit comparison
interval is $t>0$, $g(t)>0$; an actual return also requires positive
complete central outputs. As in Section~5, extend each absolute map by
zero on its own nonpassing side.

\begin{proposition}[Two-root local count]
\label{rw:finite:two-root-count}
Near every regular two-root match there are common coefficient
and section neighborhoods and a scale cutoff such that the closing
function has at most three isolated zeros, including multiplicity,
on its positive central passing interval. The corresponding
logarithmic ratio has at most two zeros and its derivative at most
one. These assertions include either sign of splitting and the
reversible match.
\end{proposition}
\begin{proof}
The exact function is
\[
 F=D_+(t)-D_-(g(t))-\rho\nu H,\qquad
 D_\pm(x)=K_\pm x^{h_\pm/\epsilon}e^{j_\pm(x)},
\]
where $h_+=p_l/\Delta+O(\epsilon)$,
$h_-=q_l/\Delta+O(\epsilon)$ and $\Delta=\sqrt{\bar a^2+4\beta}$.
All reference terms are input independent. Put
\[
 R_\pm=j_\pm'+\epsilon
       \frac{j_\pm'+xj_\pm''}{h_\pm+\epsilon xj_\pm'}.
\]
The denominator has a positive margin. Thus
$D_\pm''/D_\pm'=(h_\pm/\epsilon-1)/x+R_\pm$
and the exact cleared numerator is
\begin{equation}\label{rw:finite:two-root-numerator}
 M=tgK'=(h_--\epsilon)tg'-(h_+-\epsilon)g
        +\epsilon tg\{g'R_-(g)-R_+(t)+g''/g'\}.
\end{equation}
The multiplier $tg$ is positive on the whole passing interval.
Away from $\bar a=0$,
$h_--h_+=-\bar a/(\beta\Delta)+O(\epsilon)$ has a nonzero
margin. Differentiating \eqref{rw:finite:two-root-numerator} gives
\[
 M'=(h_--h_+)g'+(h_--\epsilon)tg''+O(\epsilon(t+g)).
\]
First choosing the section radius small, then the coefficient
neighborhood and scale small, makes $M'$ have a fixed nonzero
sign. Thus $M$ has at most one simple zero there.

Near the reversible match use the independent transverse variables
$(\bar a,\delta,\nu)$. The inverse coefficient chart follows
from the strict splitting derivative just proved; its finite jets
are bounded on common boxes. At their zero, the reflected power
coordinates may be chosen equal, giving $D_+=D_-$ and $g=t$.
The relative coefficients extend by the common linear rule of
Section~5, so the required coefficient segments are legitimate.
At zero scale the columns of the numerator are
\[
 M_{\bar a}=-\frac{t}{2\beta^{3/2}}+O(t^2),\qquad
 M_\delta=\frac1{2\beta}+O(t^2).
\]
Indeed $p_l=q_l=1/\sqrt\beta$, $\Delta=2\sqrt\beta$ and
$(h_--h_+)_{\bar a}=-1/(2\beta^{3/2})$; differentiating
$g=(t-\delta)k$ with the common reflected unit $k(0)=1$
gives the constant column. The same formulas hold after a common
positive normalization if a different reflected unit was used.
Consequently the second column has a positive margin and the
derivative of their quotient is strictly negative on a short
interval including zero. One coefficient Hadamard integral now
gives
\[
 M=\bar a A_\epsilon+\delta B_\epsilon+\nu C_\epsilon,
 \qquad B_\epsilon>c,\quad (A_\epsilon/B_\epsilon)'<-c.
\]
The finite-order preparations reserve two ordinary input derivatives
for each logarithmic rate and one coefficient derivative for this
division, which gives the displayed derivative bounds.

On the common small lower input intervals, the two large powers
and their required absolute coefficient derivatives are bounded
by $C\epsilon^{-N}e^{-b/\epsilon}$ with a fixed $b>0$.
Their exact reversible equality and the signed extensions give,
at any closing zero, $|\nu|=o(1)(|\bar a|+|\delta|)$ by
Lemma~\ref{rw:centers:relative-forcing}. If the latter two
parameters vanish, the fibre is a center and has no isolated zero.
Otherwise the strict quotient margin and relative $C^1$ forcing bound
give at most one simple zero of $M$ by
Lemma~\ref{rw:atlas:two-columns}.
Two Rolle steps give the claimed bounds for $K$ and $F$.
Both full central outputs increase with $t$, so their positivity
truncates the deficit interval to one interval. Excluded endpoints
contribute no formal closing zero.
\end{proof}

With a finite transverse upper limit, the lower possibilities are now
exhausted: both arcs are regular, the incoming arc alone loses
positivity, or both arcs lose positivity.
Lemma~\ref{rw:finite:action-order} excludes an outgoing loss alone.
The remaining finite upper limit is a saddle and is treated next.

\subsection{A finite upper saddle above regular lower arcs}
\label{rw:finite:upper-saddle-section}

Suppose $\bar m>0$ and put
\[
 p_c=\frac2{\bar c+\sqrt{\bar c^2+4\bar m}},\qquad
 q_c=\frac2{-\bar c+\sqrt{\bar c^2+4\bar m}}.
\]
We first assume $F_+>0$ on $[0,p_c]$ and $F_->0$ on
$[0,q_c]$. The perturbed separatrices meet $Y=0$ at
$P_s,Q_s$, and the returning quadrant of Section~\ref{rw:upper-resonance}
has regular outer coordinates
\[
 p=P_s-\psi_+(t),\qquad q=Q_s-\psi_-(v),
 \qquad \psi_\pm(0)=0,\quad\psi_\pm'>c>0.
\]
Retain the actual exponent $r$ of the saddle passage $v=v(t)$.
With weak references fixed independently of input,
Lemma~\ref{rw:finite:regular-lower} gives complete central maps
$s_++U(t)$, $s_-+V(v)$ and closing function
\begin{equation}\label{rw:finite:saddle-closing}
 F(t)=U(t)-V(v(t))-\rho\nu H,\qquad H>h_*>0,
 \qquad U',V'>0.
\end{equation}
The gap $H$ is independent of $t$. On a larger common domain the
regular rates
\[
 \mathsf A_+(t)=\epsilon U''/U',\qquad \mathsf A_-(v)=\epsilon V''/V'
\]
have positive margins and the finite mixed jets used below.
A physical input must traverse the chosen quadrant and both regular
outer legs, and both central outputs must be positive. These conditions
restrict a single interval; positivity of each central output is a
one-sided restriction because both maps are increasing.

The same scalar comparison will be used at secondary scale.

\begin{lemma}[Saddle logarithms and their zero count]
\label{rw:finite:saddle-counting}
Let $\epsilon,r>0$, $D=t\partial_t$, $\kappa=r-1$, and let $K$
be real analytic on $0<t<t_0$. Fix positive constants $c,C,c_1$.
$K_c,R_*,T,W$ are independent of $t$. With common bounds and
implied constants, each preparation gives, for sufficiently small
$\epsilon$, uniform section cutoffs and the stated bound with multiplicity.
\begin{enumerate}[label=\textup{(\roman*)}]
\item Suppose
\[
\begin{gathered}
 K=K_c+\epsilon\kappa\log t+\mathcal E(t),\qquad
 |K_c|\ge cR_*,\\
 |\mathcal E|+|D\mathcal E|\le CR_*\sqrt t,
 \quad R_*\ge c_1|\kappa|,\quad R_*>0.
\end{gathered}
\]
Then $Z(K)\le1$. If $\kappa=0$ there is no zero; otherwise
every zero is simple and $K'$ has sign $\kappa$ there.
\item Suppose $\kappa\ne0$ and, with the compensator of
\eqref{rw:upper:comp-map},
\[
\begin{gathered}
 DK=\epsilon\kappa+Tt-Wt\omega(t;\kappa)+R(t),\\
 W/\kappa\ge c,\quad |T|\le C|\kappa|,\quad
 |R'|\le C|\kappa|\sqrt t .
\end{gathered}
\]
First choosing $t_0$ small and then $|\kappa|$ small gives
$(DK)'/\kappa<0$, hence $Z(K)\le2$.
\item Suppose $r$ belongs to a compact positive interval with
$|r-1|\ge\delta>0$, $A,B,k_0$ have common positive upper and
lower bounds, and
\[
 (DK)'=r^2k_0At^{r-1}-B
       +O(t+t^{r-1}(t+t^r))+O(\epsilon[1+t^{r-1}]).
\]
Then $(DK)'$ is positive for $r<1$ and negative for $r>1$,
and $Z(K)\le2$.
\end{enumerate}
The bounds persist on subintervals, without using a value at $t=0$.
\end{lemma}
\begin{proof}
In \textup{(i)}, choose $t_0$ so $C\sqrt{t_0}<c/2$.
For $\kappa=0$ the constant fixes the sign. Otherwise, a zero requires
\[
 t\le e^{-cR_*/(2\epsilon|\kappa|)},\qquad
 \frac{|D\mathcal E|}{\epsilon|\kappa|}
 \le\frac{CR_*}{\epsilon|\kappa|}
          e^{-cR_*/(4\epsilon|\kappa|)}<1/2 .
\]
The last inequality is uniform since
$R_*/(\epsilon|\kappa|)\ge c_1/\epsilon$. Thus every zero is simple
and crosses with sign $\kappa$, giving at most one.

For \textup{(ii)}, put $\phi=(t\omega)'$. Then
$\phi'=-rt^{\kappa-1}<0$ and
$\phi(t_0,\kappa)\to-\log t_0-1$ as $\kappa\to0$.
Choose $t_0$ so the latter exceeds $4C/c+1$, and then choose
the ratio neighborhood so this margin persists. Throughout the interval,
\[
 (DK)'/\kappa=T/\kappa-(W/\kappa)\phi+R'/\kappa<0 .
\]
Thus $DK$ has at most one zero. In \textup{(iii)}, $t^{r-1}$
dominates for $r<1$ and tends uniformly to zero for $r>1$;
choose the section and then the scale to absorb the errors.
Rolle's inequality and $Z(K')=Z(DK)$ give both two-zero bounds.
\end{proof}

\begin{proposition}[The regular-lower upper-saddle chart]
\label{rw:finite:regular-saddle}
On common coefficient and section neighborhoods of these data, for all
sufficiently small positive scales, the actual closing function has at
most three isolated zeros, including multiplicity. This bound is uniform
through the actual ratio $r=1$. Its critical function
\begin{equation}\label{rw:finite:saddle-critical}
 K=\epsilon\log\frac{V'(v)v'}{U'}
\end{equation}
has at most two zeros on any nonidentity fibre containing a closing
zero. More generally the critical-function bound holds whenever
\eqref{rw:finite:saddle-relative-forcing} is supplied on the fixed
fibre. It holds on the whole common comparison interval $0<t<t_0$,
including inputs where a complete central output is nonpositive;
central positivity only restricts this interval when actual returns
are counted. The reversible identity fibre contributes no isolated zero.
\end{proposition}
\begin{proof}
By \ref{rw:zero-calculus} it suffices to bound $Z(K)$ by two.
Put $D=t\partial_t$ and
$E=t v''/v'-(r-1)$. The exact differentiated critical function is
\begin{equation}\label{rw:finite:saddle-M}
 M=DK=\epsilon(r-1)+\mathsf A_-(v)tv'-t\mathsf A_+(t)+\epsilon E.
\end{equation}
The outer coefficients have bounded jets through $t=v=0$.

First keep $r$ a fixed distance from one. Axis normalization as in
\eqref{rw:upper:axis-equation}, without its subsequent resonant
shears, gives $v=k_0t^r e^w$ with a bounded positive unit $k_0$.
The reciprocal integral there implies
$|D^jw|\le C\epsilon(t+t^r)$ for the finitely many input orders
needed here. Indeed
$t^r\int_t^1s^{-r}\,ds=(t-t^r)/(r-1)$; the same estimate for
its Euler derivatives follows from the integral equation and bounded
state derivatives on the rectangle. The ratio separation bounds the
denominator. Thus Lemma~\ref{rw:finite:saddle-counting}\textup{(iii)}
applies with $A=\mathsf A_-(0)$ and
$B=\mathsf A_+(0)$, proving the count on this ratio range.

Near ratio one, use the compensator preparation
\eqref{rw:upper:comp-map}, with $\kappa=r-1$ and $|\kappa|\le1/64$.
Reserve eight total orders there, so the curvature formula retains
four Euler and two coefficient orders. The actual ratio is a legal
coefficient coordinate: at the reversible face
$\partial_{\bar c}\kappa=-1/\sqrt{\bar m}+O(\epsilon)$.
Integrating the regular rates gives the exact identity
\begin{equation}\label{rw:finite:saddle-log}
\begin{aligned}
 K&=K_c+\epsilon\kappa\log t+\mathcal E(t),\qquad
 K_c=\epsilon\log\frac{rV'(0)}{U'(0)},\\
 \mathcal E(t)&=\int_0^{v(t)}\mathsf A_-(z)\,dz-\int_0^t\mathsf A_+(z)\,dz
       +\epsilon\left[w+\log\frac{r+Dw}{r}\right].
\end{aligned}
\end{equation}
The passage estimates imply
$|\mathcal E|+|D\mathcal E|\le Ct^{1/2}$.
No restriction is placed on $|\kappa\log t|$.

To compare the constant with the logarithmic term, write
$U=\alpha_U(p)e^{-A_+(p)/\epsilon}$ and
$V=\alpha_V(q)e^{-A_-(q)/\epsilon}$ here with the action notation
of Section~\ref{rw:finite:geometry}, and put
$r_U=A_+'-\epsilon(\log\alpha_U)'$, and similarly $r_V$.
These factors are positive. At the separatrix inputs one obtains
\begin{equation}\label{rw:finite:saddle-constant}
 K_c=A_+(P_s)-A_-(Q_s)
  +\epsilon\log\frac{r\alpha_V(Q_s)r_V(Q_s)\psi_-'(0)}
                          {\alpha_U(P_s)r_U(P_s)\psi_+'(0)}.
\end{equation}
The logarithmic factors have positive margins and bounded required jets.
On $\nu=0$ the regular connections give
$K_c=A_+(p_c)-A_-(q_c)+O_N(\epsilon(1+|\log\epsilon|))$.
At a limiting ratio-one point with $\bar a\ne0$ this limit is
\[
 -2\bar a\int_0^{1/\sqrt{\bar m}}
                       \frac{X^2}{F_+(X)F_-(X)}\,dX\ne0.
\]
It stays bounded away from zero in a smaller neighborhood.

The coordinates of Section~\ref{rw:upper-resonance} satisfy
$v=t$ and $U=V$ on $\bar a=\kappa=\nu=0$. Choose the coefficient
contractions within the same quadrant and enlarged regular lower
tubes, so these exact identities can be used for division. With
$n=|\bar a|+|\kappa|+|\nu|$, division along these contractions gives
\begin{equation}\label{rw:finite:saddle-center-estimates}
 |\mathcal E|+|D\mathcal E|\le Cnt^{1/2},\qquad
 K_c=\bar a A_c+\kappa B_c+\nu C_c,
 \quad A_c<-c,\quad |B_c|+|C_c|\le C.
\end{equation}
To see the strict sign, differentiate \eqref{rw:finite:saddle-constant}
at fixed actual $\kappa$; there
$\partial_{\bar a}\bar c=O(\epsilon)$ and
\[
 A_c\longrightarrow-2\int_0^{1/\sqrt{\bar m}}
                         \frac{X^2}{(1+\bar BX^2)^2}\,dX<0.
\]
Coefficient derivatives of $t^r$ only add fixed powers of $|\log t|$,
absorbed by the gap between $t^{1-1/64}$ and $t^{1/2}$.

On these regular lower tubes the actions have a common positive
lower bound $S_*>0$. The absolute quantities $U,V$ and their
finite coefficient derivatives are bounded by
$C\epsilon^{-M}e^{-S_*/\epsilon}$; composition with
$t^r$ and its coefficient derivatives preserves this bound uniformly
on the section. The center identity and coefficient division therefore
give the hypothesis of Lemma~\ref{rw:centers:relative-forcing} with
$I=|\bar a|+|\kappa|$. Every fibre containing a closing zero satisfies
\begin{equation}\label{rw:finite:saddle-relative-forcing}
 |\nu|\le\eta(\epsilon)(|\bar a|+|\kappa|),
 \qquad\eta(\epsilon)\longrightarrow0.
\end{equation}
If both generators vanish, this forces the exact identity fibre.

For a nonzero limiting $\bar a$, the constant and remainder estimates
satisfy Lemma~\ref{rw:finite:saddle-counting}\textup{(i)} with $R_*=1$.
Near the reversible face fix a large $K_0$. In the cone
$|\bar a|>K_0|\kappa|$, the divided constant and relative forcing
give the same hypotheses with $R_*=|\bar a|$.

In the remaining cone $|\bar a|\le K_0|\kappa|$, the compensator
controls the derivative. A nonidentity fibre has $\kappa\ne0$ and
$n\le C_{K_0}|\kappa|$.
Set $A=\mathsf A_-(0)$ and $B_0=\mathsf A_+(0)$. Using
$t^r=t-\kappa t\omega(t;\kappa)$ in
\eqref{rw:finite:saddle-M} and the curvature formula gives
\begin{equation}\label{rw:finite:saddle-compensator}
\begin{aligned}
 M&=\epsilon\kappa+Tt-Wt\omega+R,\\
 T&=rA-B_0+\epsilon\sigma(2+1/r),\qquad
 W=r(A\kappa+2\epsilon\sigma).
\end{aligned}
\end{equation}
The prepared coefficients vanish at zero scale and on the reversible
face, so $|\sigma|\le C\epsilon n$.
All remaining terms vanish on that face as well. The reserved mixed
jets of the passage and the regular rates, followed by one coefficient
division, give $|D^jR|\le Cnt^{3/2}$, $j\le2$.
Explicitly, the undivided products have powers at least
$t^{2-2/64}$ times fixed logarithmic powers; their coefficient
variations retain a power strictly larger than $3/2$.
Also $|T|\le Cn$, and the positive margin of $A$ gives
$W/\kappa\ge c>0$ for small scale. Hence
$|R'|\le C_{K_0}|\kappa|\sqrt t$.
Lemma~\ref{rw:finite:saddle-counting}\textup{(ii)} now gives
$Z(K)\le2$ in the remaining cone. Together with the separated-ratio
case this proves the critical-function bound; \ref{rw:zero-calculus}
gives $Z(F)\le3$ on the physical subinterval.
\end{proof}

\begin{lemma}[Coefficient and endpoint signs]
\label{rw:finite:saddle-entry-signs}
On a regular-lower saddle fibre with
\eqref{rw:finite:saddle-relative-forcing}, suppose an incoming regular
overlap supplies
\[
 \tau\sigma_1<0,\quad |\tau|\le C|\sigma_1|,
 \qquad\operatorname{sign}K=-\operatorname{sign}\sigma_1,
 \qquad \tau=\bar a+\bar c,\quad
 \sigma_1=\bar c(\bar B+\bar m).
\]
Then $\bar a\bar c<0$, $|\bar a|\le C_1|\bar c|$,
$\operatorname{sign}\kappa=-\operatorname{sign}\bar c$, and
$K(0+)$ has sign $\bar c$, opposite to the supplied overlap sign.
The approach toward this saddle is increasing $p$, hence decreasing
$t$. The endpoint limit refers to the same relative critical function
on its common extension $0<t<t_0$; if a central-output condition cuts
off that interval, it does not assert a positive central return at
the excluded endpoint. These are entrance data for the interval
argument of Section~12.
\end{lemma}
\begin{proof}
Regular lower positivity gives the identity and inequality
\[
 0< F_+(p_c)/p_c+F_-(q_c)/q_c
                  =(\bar B+\bar m)(p_c+q_c).
\]
Thus $\sigma_1$ has the sign of $\bar c$. The asserted coefficient
cone follows from the two hypotheses on $\tau$.
For the actual saddle, $\kappa=-\operatorname{tr}/\lambda_u$.
The zero-scale divided field has trace $\bar c+\nu$, so near the
reversible face the expansion is
\[
 \kappa=-\bar c/\sqrt{\bar m}
 +O(\bar c^2+\epsilon|\bar c|+\epsilon|\bar a|+|\nu|).
\]
The coefficient cone and the relative forcing absorb the error,
first giving $|\kappa|\le C|\bar c|$ and then its opposite sign.
Away from $\bar c=0$ the limiting formula
$\kappa_0=-2\bar c/(\sqrt{\bar c^2+4\bar m}+\bar c)$
has a fixed nonzero margin, preserved at small scale and forcing.
Finally \eqref{rw:finite:saddle-log} has bounded remainder and
$\epsilon\kappa\log t$ has sign $-\kappa$ as $t\downarrow0$.
The coefficient cone and the supplied overlap sign therefore give
the required endpoint sign on the comparison continuation.
\end{proof}

\subsection{An upper saddle with one incoming lower loss}
\label{rw:finite:upper-loss-section}

Keep the outgoing lower arc regular and let the incoming arc lose
positivity. A simple first incoming root must coincide with $q_c$;
a double root can also occur before it. In both cases use the
whole-interval nonnegativity required by
Section~\ref{rw:finite:geometry}, rather than only the value at $q_c$.

\begin{proposition}[The incoming-loss upper-saddle chart]
\label{rw:finite:upper-loss}
On common neighborhoods of either such configuration, the actual
closing function has at most three isolated zeros, with multiplicity.
For the actual ratio $r\ge1$, its critical function is strictly
negative and the closing function is strictly increasing, so its bound
is one. For $r<1$ the multiplied critical numerator below is strictly
increasing, giving at most two critical zeros and three closing zeros.
These assertions hold on the entire actual passing interval, even
when its boundary moves away from the upper separatrix input.
\end{proposition}
\begin{proof}
For a simple incoming root, Proposition~\ref{rw:lower-preparation} gives
$V(v)=D_-(g(v))$, where $g'>c>0$ and passage requires $g(v)>0$.
Write $\lambda=h_-/\epsilon$, $L=\epsilon\lambda$ and
\[
 \frac{D_-''(g)}{D_-'(g)}=\frac{\lambda-1}{g}+R_-(g),
 \qquad C(v)=R_-(g(v))g'(v)+g''(v)/g'(v).
\]
The remainder and its needed ordinary jets are bounded, by the
power preparation and the positive large-exponent factor in $D_-'$.
The critical function remains
$K=\epsilon\log(V'(v)v'/U')$. Its exact numerator, multiplied
by the positive factor $tg$, is
\begin{equation}\label{rw:finite:upper-simple-numerator}
 M_s=(L-\epsilon)tg'v'+\epsilon g(\kappa+E)
                 +\epsilon tgv'C-\mathsf A_+(t)tg.
\end{equation}

For a double root, use the increasing Morse slow entry $g(v)$,
$g'>c>0$, and write $f(v)=\mu+g(v)^2$. If the root precedes
$q_c$, choose one phase interval containing both the minimum and
the compact upper-entry range. The bound $F_I\!{}^{\prime\prime}>1$
and the positive finite horizontal coordinate
give one increasing Morse coordinate and one bounded positive time
unit on that entire interval. Fixed regular pieces may be included
in the connections. The moving-entry relative formula,
including any regular normal input as part of its relative amplitude,
gives
\[
 \frac{V''(v)}{V'(v)}
 =\frac{\mathsf A_-(g(v))g'(v)}{\epsilon f(v)}+\frac{g''(v)}{g'(v)}.
\]
Here $\mathsf A_-$ has a positive margin and bounded jets by
Lemma~\ref{rw:transitions:moving-rate}. Put
$J=\mathsf A_-(g)g'$ and $C=g''/g'$. The exact numerator, now with
positive factor $tf$, is
\begin{equation}\label{rw:finite:upper-double-numerator}
 M_d=Jtv'+\epsilon f(\kappa+E)+\epsilon ftv'C-\mathsf A_+(t)tf.
\end{equation}
Passage requires $\mu+s^2>0$ at every point between $g(v)$ and
the fixed exit. Since $g$ and $v$ increase, the permitted inputs form
an interval. The physical interval is its intersection with the upper
quadrant, regular connection domains and positive central-output
conditions. An excluded $t=0$ supplies no physical endpoint value.

For $r\ge1$ choose a fixed incoming section $R$ before its limiting
first root such that $A_-(R)>A_{\rm out}+3$. The incoming action
diverges at either a simple or double root, whereas the outgoing
action stays bounded. Choose a positive section gap and shrink the
neighborhood so that every incoming endpoint lies beyond $R$ by
that gap. The upper derivative
$v'=t^{r-1}e^w(r+Dw)$ is uniformly bounded for $r\ge1$, including
$r=1$; the regular outer derivatives are bounded too.
Lemma~\ref{rw:finite:incoming-contraction} therefore gives
$V'(v)v'<U'(t)$ on the whole passing interval, hence $K<0$.
It applies to the actual incoming arc and its reference difference;
no derivative estimate on its long transit time is required.

For $r<1$ separated from one, the first term in the derivative of
either numerator is at least $ct^{r-1}$. The other terms are
$O(1)+o(t^{r-1})+O(\epsilon t^{r-1})$ by the reciprocal upper
estimate used above. Choosing section and then scale proves $M'>0$.
As $r\uparrow1$, this argument needs the sign of the resonant
curvature term. We compute that sign in the original field.
At a limiting ratio-one configuration, $\bar c=0$ and $p_c=q_c$.
Incoming loss with outgoing positivity implies $\bar a>0$.
For a simple loss this follows from
$F_+(p_c)-F_-(p_c)=2\bar ap_c>0$; for a double loss its positive
root has $\bar a=2\sqrt{\bar B}>0$.
Appendix~\ref{rw:app:upper-resonant-sign} derives from the original
field and its exact zero-trace equation
\[
 \sigma|_{\kappa=\nu=0}
  =-\chi_{\rm loc}^2\epsilon\bar a
                       [\bar m^{5/2}+O(\epsilon)].
\]
Consequently $\sigma<-c\epsilon$ on fixed smaller ratio and
coefficient neighborhoods of this loss configuration.
The actual curvature formula is therefore
$E=-\sigma Q+R_E$, where $Q,Q'>0$ and
$|D^jR_E|\le C\epsilon t^{7/4}$ for the orders used here.

For the simple numerator, group its differentiated terms as
\[
 M_s'=(rL-\epsilon)g'v'-\mathsf A_+(0)g(0)
       +Lg'v'E+\epsilon gE'+\mathcal R_s,
\]
where
$|\mathcal R_s|\le Ct^{r-1}(v+t+\epsilon)$ on the smaller
near-resonant section. This bound follows directly by differentiating
\eqref{rw:finite:upper-simple-numerator}: the additional quadratic
velocity terms contain $tv'=O(v)$, and $g(v)-g(0)=O(v)$.
Terms with $C$ carry $\epsilon$; terms with its extra factor $E$
use $|E|\le C\epsilon t^{1-1/64}(1+|\log t|)$.
The first term is at least $ct^{r-1}$, whereas $g(0)$ tends to zero
at the limiting simple loss. The principal parts of the two displayed
curvature terms are nonnegative because $g>0$ on passage. Their
remainders are bounded by $C\epsilon t^{3/4}(1+v')$ and are
absorbed by the positive margin. Thus $M_s'>0$.

For the double numerator the same differentiation gives
\[
 M_d'=r^2J(0)t^{r-1}-\mathsf A_+(0)f(0)
 +(J+\epsilon f_v)Ev'+\epsilon fE'+\mathcal R_d,
 \qquad |\mathcal R_d|\le Ct^{r-1}(v+t+\epsilon).
\]
Here $f(0)$ need not tend to zero when the double root precedes the
upper endpoint, so compare the two constant terms together.
At zero scale and ratio one the incoming phase is
$g=1/q-\bar a/2$, $\mu=\bar B-\bar a^2/4$.
Put $k_0=g'(0)>0$. The symmetric upper connections then have
$p_t(0)=q_v(0)=-k_0p_c^2$. Division of the lower slow time by
$x=1/q$ makes its limiting normal rate $1/x=q$.
Consequently
\[
 J(0)=p_ck_0,\qquad \mathsf A_+(0)=\frac{k_0p_c^3}{F_+(p_c)},
 \qquad f(0)=\frac{F_-(p_c)}{p_c^2}.
\]
Their exact limiting difference is
\begin{equation}\label{rw:finite:upper-double-gap}
 J(0)-\mathsf A_+(0)f(0)
       =\frac{2\bar a k_0p_c^2}{F_+(p_c)}>0.
\end{equation}
For nearby $r<1$, the factor $t^{r-1}\ge1$ preserves a fixed
positive margin after reducing the ratio neighborhood; if $f(0)<0$
the estimate improves. Thus this argument includes a double root
before the limiting upper endpoint. Since $J+\epsilon f_v>c$
and $f>0$ on the actual interval, the principal curvature terms are
again nonnegative. Their remainder and $\mathcal R_d$ are smaller
than this margin after fixed section and scale restrictions. Hence
$M_d'>0$ on the entire passing interval.

In both cases $M$ is a positive multiple of $K'$ and is strictly
increasing. Thus \ref{rw:zero-calculus} gives at most two critical
zeros and three closing zeros on the entire physical interval, including
multiplicities and analytic endpoint germs. The entrance argument in
Section~12 joins this interval to its preceding regular part.
\end{proof}

\subsection{Two lower roots at a finite upper saddle}
\label{rw:finite:three-saddle}

When both simple lower roots meet an upper saddle, all three
passages can lose their exits together. Fix $\bar B=-\beta$,
$0<\beta_*\le\beta\le\beta^*$, near the triangle center
$\bar m=\beta$, $\bar c=-\bar a$, $\nu=0$. The actual upper
ratio $r$ stays in a fixed compact subset of $(0,\infty)$. Retaining
both lower deficits as well as this ratio allows all three passages
to degenerate independently. Use the actual splitting coordinates of Lemma~\ref{rw:coordinates:finite-splittings}, and the
upper coordinates and regular legs of
Lemma~\ref{rw:coordinates:two-center-upper} near the reversible
intersection. On the remaining compact parts, extend just the triangle
coordinates $L_\pm L_0^{\epsilon\beta}$ from
Lemma~\ref{rw:centers:triangle}, and straighten and normalize the
neighboring axes by the same construction. The axis graph denominators
have positive margins on every compact positive-ratio interval; no
resonant shear is needed for this extension. Its orbital remainder is
$\epsilon$ times a bounded function and vanishes on the triangle
component. The same regular outer legs therefore give
\[
 u=G_+(t)=u_0+tA_0(t),\qquad
 z=G_-(v)=z_0+vB_0(v),\qquad v=t^re^{w(t)},
\]
where $A_0,B_0,G_+',G_-'>c>0$. The lower maps are those of
Proposition~\ref{rw:lower-preparation}:
\begin{equation}\label{rw:finite:three-powers}
 D_\pm(s)=K_\pm s^{h_\pm/\epsilon}e^{j_\pm(s)},\qquad
 K_\pm>0,\quad h_\pm>h_*>0,\quad j_\pm(0)=0.
\end{equation}
The upper passage and both lower deficit passages define the signed
central comparison on
\begin{equation}\label{rw:finite:three-domain}
 0<t<t_0,\qquad u>0,\qquad z>0.
\end{equation}
Since $u,z$ increase, this is an interval. The closing function there is
\[
 F(t)=D_+(u)-D_-(z)-\sqrt\epsilon\,\nu H,\qquad H>h_*>0.
\]
The weak references, and hence $H$, are independent of $t$.
The complete central outputs are $s_++D_+(u)$ and $s_-+D_-(z)$.
Their positivity is an additional requirement for an actual return,
beyond the deficit conditions in \eqref{rw:finite:three-domain}.
The Gaussian reference calculation in
Lemma~\ref{rw:finite:regular-lower} applies on the fixed inner cores:
$F_\pm$ have positive margins there even when the outer roots are
lost. It gives $H=\sqrt{2\pi}+O(\sqrt\epsilon)$ and, more precisely,
\[
 s_+=-\sqrt\epsilon\,\nu H_+,
 \qquad s_-=\sqrt\epsilon\,\nu H_-,
 \qquad H_+,H_->c>0,\qquad H=H_++H_-.
\]
At a zero of $F$ the two outputs agree. For $\nu\ge0$ the
incoming one is positive, and for $\nu\le0$ the outgoing one is
positive, since $D_\pm>0$. Thus their common output is always
positive at a closing zero. In each core the original height
$\mathfrak z=1+y$ has derivative of fixed sign on $\mathfrak z=0$,
namely the sign of $\nu$ (with the reflected sign on the other
core). An orbit with positive entrance and exit cannot cross that
line and return. Together with the strict deficit passages and finite
transverse connections, this proves that each closing zero gives an actual
positive-height return. At nonclosing inputs we continue to use the
signed comparison.

The numerator calculation below uses coefficient division as well as
input differentiation. Reserve twelve total finite orders in the lower
and upper preparations before shrinking their construction domains.
This supplies six ordinary input and two coefficient derivatives of the regular connectors and
lower relative coefficients, four Euler input and four coefficient
derivatives of the upper relative logarithm, and the two coefficient
derivatives spent in product division below. The lower graph can be
constructed at order $23$ by the $N+11$ reserve in
Proposition~\ref{rw:lower-preparation}. Scale derivatives are Euler
derivatives. The ordinary input derivatives of the lower connectors
are not asserted for the singular composite $t^r$ at zero.

Use the common signed relative extension and absolute zero extension
after Proposition~\ref{rw:generator-division} when a coefficient
segment changes a deficit's sign.
For any fixed desired $S_*>0$, first reduce the input radius $d<1$
so that a half of $d^{h_*/\epsilon}$ absorbs the upper bound for
$\exp(\epsilon\log K_\pm/\epsilon)$, and then reduce $\epsilon$
so that the remaining half-power absorbs all reserved derivatives.
This gives
\begin{equation}\label{rw:finite:three-flat}
 |\partial_s^i\partial_{\mathbf b}^{\alpha}\widehat D_\pm|
       \le C\epsilon^{-M}e^{-S_*/\epsilon}
\end{equation}
through the required finite orders. These extensions preserve both
center identities by that construction.
Keep the original $t$ fixed on the coefficient segments and impose
central positivity on actual returns.

Put
\[
 \mathcal K=\epsilon\log
 \frac{D_-'(z)G_-'(v)v'}{D_+'(u)G_+'(t)},\qquad
 D=t\partial_t,\qquad \kappa=r-1,\qquad
 E=\frac{tv''}{v'}-\kappa.
\]
Ordinary primes on $G_\pm$ and $D_\pm$ mean their own input
derivatives. For use in the chain rule define
\[
 R_\pm(s)=j_\pm'(s)+
       \frac{j_\pm'(s)+sj_\pm''(s)}{h_\pm/\epsilon+sj_\pm'(s)}.
\]
Then $D_\pm''/D_\pm'=(h_\pm/\epsilon-1)/s+R_\pm$ on $s>0$.
The displayed formula also defines bounded auxiliary coefficients on
the signed interval; no logarithm of a zero extension is taken.

\begin{proposition}[Complete three-saddle numerator]
\label{rw:finite:three-numerator}
On the comparison interval the following function is exactly
$tuz\,\mathcal K'$:
\begin{equation}\label{rw:finite:three-M}
\begin{aligned}
 M={}&(h_--\epsilon)uDz-(h_+-\epsilon)zDu
              +\epsilon uz(\kappa+E)\\
 &+\epsilon uz\left[
 R_-(z)Dz-R_+(u)Du+
 \frac{G_-''(v)}{G_-'(v)}Dv-
 \frac{G_+''(t)}{G_+'(t)}t\right].
\end{aligned}
\end{equation}
This formula defines its auxiliary extension. Its endpoint value is
\[
 M(0+)=\epsilon\kappa u_0z_0.
\]
Near ratio one put
$\xi_1=(u_0+z_0)/2$, $\xi_2=(u_0-z_0)/2$ and
$I=|\xi_2|+|\kappa\xi_1|$. Then
\begin{equation}\label{rw:finite:three-comp}
\begin{aligned}
 M&=C_0+Tt-Wt\omega(t;\kappa)+\mathcal R_3,\\
 C_0&=\epsilon\kappa u_0z_0,\\
 P&=u_0z_1(rh_--\epsilon)+\epsilon u_0z_0rJ_-,\\
 N&=z_0u_1(h_+-\epsilon r)+\epsilon u_0z_0J_+,\\
 T&=P-N+\epsilon u_0z_0\sigma_s(2+1/r),\\
 W&=\kappa P+2\epsilon u_0z_0\sigma_s r,
\end{aligned}
\end{equation}
where $u_1=G_+'(0)$, $z_1=G_-'(0)$ and
\[
 J_+=u_1R_+(u_0)+G_+''(0)/u_1,\qquad
 J_-=z_1R_-(z_0)+G_-''(0)/z_1.
\]
Here $\sigma_s$ denotes the resonant coefficient $\sigma$ of
\eqref{rw:upper:comp-map} in the present center-preserving coordinates.
For $|\kappa|\le1/64$,
\begin{equation}\label{rw:finite:three-remainder}
 |D^j\mathcal R_3|\le C(I+|\nu|)t^{7/4},\qquad 0\le j\le2.
\end{equation}
If this fibre has one actual closing zero, then
\begin{equation}\label{rw:finite:three-forcing}
 |\nu|\le\delta_\epsilon I,\qquad
 \delta_\epsilon=C\epsilon^{-M-1/2}e^{-S_*/\epsilon}\longrightarrow0.
\end{equation}
\end{proposition}

\begin{proof}
Differentiate $\mathcal K$ before clearing its denominators.
The two power rates give the first two terms in
\eqref{rw:finite:three-M}; the upper second derivative gives
$\epsilon uz(\kappa+E)$. The ordinary curvature of each regular leg
gives the last two terms in its bracket. Thus the numerator retains
every coordinate contribution to the complete central derivative.

To verify the expansion, first regard $t$ and $t^r$ as independent
small quantities. The endpoint Taylor formulas give
\[
 u=u_0+u_1t+O(t^2),\qquad
 z=z_0+z_1t^r+\text{terms of order greater than one}.
\]
The terms linear in these two quantities are
\[
 C_0+Pt^r-Nt-\epsilon u_0z_0\sigma_s Q(t),
 \qquad Q=2rt\omega-(2+1/r)t.
\]
In particular the $J_\pm$ contributions are present even when the
connectors are nonlinear. The exact equality
$t^r=t-\kappa t\omega$ gives \eqref{rw:finite:three-comp}.

We first estimate the remainder uniformly on the auxiliary box,
before dividing by the center generators. Write $\mathcal F_a$ for a
function whose selected
$D^j\partial_{\mathbf b}^{\alpha}$ derivatives are $O(t^a)$,
$j\le2$, $|\alpha|\le2$. With $a=15/16$ the functions
$t,t^r,t\omega$ belong to $\mathcal F_a$ on
$|\kappa|\le1/64$; every fixed parameter logarithm is absorbed by
$63/64-a>0$. The upper preparation gives
$v-t^r,Dv-rt^r\in\mathcal F_{2a}$ and
$E+\sigma_sQ\in\mathcal F_{7/4}$.
Integral Taylor remainders of $G_\pm$, $R_\pm$ and their regular
curvatures then show that every omitted product has at least two
$\mathcal F_a$ factors or the displayed curvature remainder.
As $2a=15/8>7/4$, it follows that
$D^j\partial_{\mathbf b}^{\alpha}\mathcal R_3=O(t^{7/4})$ on
the entire common auxiliary box.

The triangle center gives vanishing of $M$ and its retained
coefficients on $\xi_1=\xi_2=\nu=0$; the reflected extension
identities give the same vanishing on $\xi_2=\kappa=\nu=0$.
These two axes explain the generator pair $(\xi_2,\kappa\xi_1)$.
Apply the product-ideal division of
Proposition~\ref{rw:generator-division} to $\mathcal R_3$.
Its explicit mixed integral in $\xi_1,\kappa$ spends two coefficient
derivatives and preserves $t^{7/4}$, proving
\eqref{rw:finite:three-remainder}.

Apply the identical division to
$\widehat D_+(G_+(t))-\widehat D_-(G_-(v(t)))$.
Equation~\eqref{rw:finite:three-flat} bounds its divided coefficients
by $C\epsilon^{-M}e^{-S_*/\epsilon}$. At a physical closing zero
Lemma~\ref{rw:centers:relative-forcing} now proves
\eqref{rw:finite:three-forcing}. This is a coefficient constraint
on the whole fibre. If $I=0$, it forces $\nu=0$ and the exact center
identity contributes no isolated zero.
\end{proof}

\paragraph{The scalar compensator calculation.}
The same scalar calculation controls the derivative at a single upper
saddle and the numerator here. Keeping $t\omega$ intact makes its
bounds uniform even when $|\kappa\log t|$ is unbounded.

\begin{lemma}\label{rw:finite:compensator-calculus}
For $|\kappa|\le1/64$ set $r=1+\kappa$ and
\[
 M=C+Tt-Wt\omega(t;\kappa)+R,\qquad L=|T|+|W|>0.
\]
Assume $|R'|\le C_*Lt^{1/2}$ and
$|R''|\le C_*Lt^{-1/2}$. On a common sufficiently small
$(0,t_0)$, $M$ has at most two zeros, with multiplicity.
If it has two, then
\begin{equation}\label{rw:finite:compensator-signs}
 TW>0,\qquad (C+R(0))W>0.
\end{equation}
There is also the following one-zero version. For fixed positive
$c_*,C_*$, if
\[
 W/\kappa\ge c_*,\quad |T|\le C_*|\kappa|,\quad
 |R'|\le C_*|\kappa|t^{1/2},\quad \kappa\ne0,
\]
then first choosing $t_0$ small and then $|\kappa|$ small makes
$M'$ have sign $-\operatorname{sign}\kappa$ throughout the interval.
The constant $C$ is unrestricted in both assertions.
\end{lemma}

\begin{proof}
Let $\phi=(t\omega)'$. Direct differentiation gives
$\phi'=-rt^{\kappa-1}$. On $0<t<e^{-6}$,
$1\le\phi\le16t^{-1/8}$. At any zero of $M'$,
\[
 T=W\phi-R',\qquad
 |W|\ge\frac{L}{2(1+\phi)}
\]
once $C_*t_0^{1/2}<1/2$. Thus $W\ne0$ at such a zero and
\[
 \frac{|R''|}{|W|rt^{\kappa-1}}
 \le \frac{2C_*}{r}(1+\phi)t^{1/2-\kappa}
 \le C C_*t^{23/64}<1/2
\]
after a second common reduction. Every zero of $M'$ is therefore
simple and crosses with sign $W$. There can be at most one such
crossing; Rolle's theorem proves the two-zero bound, including
multiplicity. If $W=0$, the estimate on $R'$ excludes a zero of $M'$.

When $TW\le0$, $|T-W\phi|\ge L>|R'|$, so two zeros are impossible.
If $W>0$, the unique possible critical point is a minimum. Two zeros
force strict decrease to the first of them and hence $M(0+)>0$.
For $W<0$ all signs reverse. The integrable bound on $R'$ supplies
$R(0)$ and proves \eqref{rw:finite:compensator-signs}.
The last assertion is the $\phi$ estimate in the proof of
Lemma~\ref{rw:finite:saddle-counting}\textup{(ii)}, applied to
$M'=T-W\phi+R'$.
\end{proof}

\paragraph{The coefficient columns near ratio one.}
To apply the scalar lemma, we must control its coefficient norm
$|T|+|W|$ by the center generators. Work on the same inverse box
as above. Transversality of $T$ to the reversible component and an
exact factorization of $W$ give the required determinant margin.
At the reversible intersection, let $k_*=G_+'(0)=G_-'(0)>0$.
There $h_+=h_-=1/(2\beta)$ exactly. Put
\[
 A_c=k_*(1/(2\beta)-\epsilon)>c>0.
\]
The coefficient $T$ has $\partial_{\xi_2}T=2A_c$ there.
Moreover the functions $\tau_1$, $\sigma_1$ and $\sigma_s/\epsilon$
vanish on both center axes, so their divisions have bounded columns:
\begin{equation}\label{rw:finite:three-directed-generators}
\begin{aligned}
 \tau_1&=a_\tau\xi_2+b_\tau\kappa\xi_1+c_\tau\nu,
                         &a_\tau&>c>0,\\
 \sigma_s/\epsilon&=-a_s\xi_2+b_s\kappa\xi_1+c_s\nu,
                         &a_s&>c>0,\\
 |\tau_1|+|\sigma_1|&\le C(I+|\nu|).
\end{aligned}
\end{equation}
The positive $a_\tau$ is the inverse of the positive transverse
column in \eqref{rw:coordinates:finite-jacobian}, continued with
its coefficient derivative margin. We verify the second direction
from the original upper field.

The coordinates of Lemma~\ref{rw:coordinates:two-center-upper}
are affine at zero scale, so the first-order calculation in
Appendix~\ref{rw:app:upper-resonant-sign} applies. Absorbing their
positive linear section factors in $\chi_{\rm loc}^2$, it gives on
$\kappa=\nu=0$
\[
 \sigma_s=-\chi_{\rm loc}^2\epsilon\bar a
                  [\bar m^{5/2}+O(\epsilon)],\qquad\chi_{\rm loc}>c.
\]
The factors $\epsilon$ and $\bar a$ are exact, by the affine
zero-scale field and the reversible identity. The finite coordinate
Jacobian gives $\partial_{\bar a}\xi_2>c$ at fixed $\kappa$.
Hence the $\xi_2$ column is strictly negative at the intersection.
The integral division, with the uniform second coefficient jets,
retains that sign on a smaller box and proves
\eqref{rw:finite:three-directed-generators}.

The $W$ column has an exact positive factor on the finite-scale
coefficient box. Define
\[
 Q_*=z_1(rh_--\epsilon)+\epsilon z_0rJ_->c>0,\qquad
 Q^*=u_1(h_+-\epsilon r)+\epsilon u_0J_+>c>0.
\]
Exactly, $P=u_0Q_*$ and $N=z_0Q^*$. On $\xi_2=\nu=0$ write
$\sigma_s=\epsilon b_s\kappa\xi_1$. Then
\begin{equation}\label{rw:finite:three-W-factor}
 W=\kappa\xi_1[Q_*+2\epsilon^2rb_s\xi_1^2].
\end{equation}
The bracket has a uniform positive margin.
Remove $\xi_2,\nu$ by single integral division. On $\nu=0$ this gives
\[
 \binom TW=
 \begin{pmatrix}a_T&b_T\\a_W&b_W\end{pmatrix}
 \binom{\xi_2}{\kappa\xi_1},
 \qquad a_T>c,\quad b_W>c,\quad |b_T|\le C.
\]
The exact formula for $W$ also gives
$|a_W|\le C|\kappa|+C\epsilon^2(|\xi_1|+|\xi_2|)^2$.
Consequently the determinant has a positive margin after reducing
the box; restoring $\nu$ and using
\eqref{rw:finite:three-forcing} yields
\begin{equation}\label{rw:finite:three-column-norm}
 cI\le |T|+|W|\le CI
\end{equation}
on every noncenter fibre with a closing zero. The division uses the
coordinates $\kappa,\xi_1$ themselves; their product
$\kappa\xi_1$ is a generator.

\paragraph{Critical points on the deficit domain.}
The deficit inequalities now constrain a critical point of $M$.
This step uses $u,z>0$ and therefore applies on the comparison
interval, rather than on the whole auxiliary box.

\begin{lemma}\label{rw:finite:three-critical-gaps}
Near ratio one, under \eqref{rw:finite:three-forcing} on a noncenter
fibre, every zero of $M'$ in the comparison interval has $u_0z_0>0$.
\end{lemma}

\begin{proof}
By \eqref{rw:finite:three-comp} and
$1-\kappa(t\omega)'=rt^\kappa$,
\[
 M'=rP t^\kappa-N-\epsilon u_0z_0\sigma_sQ'
                                 +\mathcal R_3'.
\]
Suppose $u_0z_0\le0$, and put
$S(t)=|u_0|t^\kappa+|z_0|$. The principal terms have the same
sign and modulus at least $cS(t)$; the noncenter assumption and
forcing exclude $S=0$.
If $u_0<0$, the condition $u>0$ gives $|u_0|\le Ct$.
If $z_0<0$, $z>0$ gives $|z_0|\le Cv\le Ct^r$.
In either case
\[
 |u_0z_0|\le CtS(t).
\]
Since $|Q'|\le Ct^{-1/8}$, its curvature contribution is bounded
by $C\epsilon|\sigma_s|t^{7/8}S(t)$.
Also $I\le C(|u_0|+|z_0|)$ and
$S(t)\ge t^{1/64}(|u_0|+|z_0|)$, whence
\[
 |\mathcal R_3'|\le CIt^{3/4}
                      \le Ct^{47/64}S(t).
\]
A common small $t_0$ makes both errors less than half the principal
margin. Therefore $M'$ cannot vanish.
\end{proof}

\paragraph{Ratios separated from one.}
Away from ratio one, a power comparison replaces the compensator.
We must also recover the signs of the original coefficients: the
curvature crossing alone does not determine them.
After reflection it is enough to fix $r\ge1+\delta$, $\delta>0$,
on a compact ratio interval. Put $I_a=|u_0|+|z_0|$.
The triangle-preserving coordinates give the unsheared equation
\[
 V_U=-rV/U-V^2R(U,V),\qquad
 R=\epsilon(u_0R_u+z_0R_z+\nu R_\nu),
\]
with bounded finite jets. Its exact reciprocal integral yields
\[
 v=t^re^w,\quad |D^jw|\le C\epsilon(I_a+|\nu|)t,\quad
 v'\asymp t^{r-1},\quad v''\ge c\,v'/t.
\]
Indeed $t^r\int_t^1s^{-r}\,\dd s\le Ct$, and its finitely many
Euler derivatives have the same bound. The differentiated integral
equations have small integral coefficients; iteration retains these
bounds. Fixed coefficient logarithms are absorbed with a portion of
the positive margin $\delta$. No near-resonant shear is used here.

Single-center division of the absolute closing difference gives
$|\nu|\le\delta_\epsilon I_a$ on every fibre with a closing zero.
Direct integral division of \eqref{rw:finite:three-M} then gives,
on the full positive $t$ interval of the auxiliary extension,
\begin{equation}\label{rw:finite:three-away-derivative}
 M'=u_0v'\mathcal A-z_0\mathcal B+\nu v\mathcal C,
\end{equation}
where
\[
 \mathcal A,\mathcal B>c>0,\quad
 |\mathcal A'|\le C,\quad
 |\mathcal B'|\le C(1+|u_0|/t),\quad
 |(v\mathcal C)'|\le Cv/t.
\]
These bounds follow by separating the two unmixed deficit terms
from their mixed connector term. On $u_0=z_0=0$, the numerator is
$tv\nu$ times a bounded coefficient, by division from the exact
triangle identity.
On $z_0=0$, divide $M$ by $v$ and then by $u_0$; its endpoint
coefficient is $(rh_--\epsilon)z_1>c$.
On $u_0=0$, divide by $t$ and then by $z_0$; its endpoint
coefficient is $-(h_+-\epsilon r)u_1<-c$.
The remaining mixed term has its constant
$\epsilon\kappa u_0z_0$ explicitly removed. Its first Euler
derivatives are $O(t)$ and its second ordinary derivative is
$O(1/t)$. The bounded connector derivatives and the preceding
bounds on $w$ prove these estimates term by term in
\eqref{rw:finite:three-M}; integral division in $u_0,z_0$
preserves them. Differentiating the two unmixed terms gives
\eqref{rw:finite:three-away-derivative}; the mixed second derivative
accounts for $|u_0|/t$ in $\mathcal B'$. Ratios of $v$ at points
of a coefficient integration segment are positive bounded units,
since $r$ is held fixed and the corresponding difference of $w$
is $O(\epsilon(I_a+|\nu|)t)$.

At a zero of $M'$, its two principal terms must balance. The
$\nu v$ term is smaller than their sum since
$v\le Ctv'$ and $v\le Ct$. Thus
\begin{equation}\label{rw:finite:three-away-balance}
 u_0z_0>0,\qquad |z_0|\asymp|u_0|v'
                      \asymp |u_0|t^{r-1}.
\end{equation}
At such a point,
$M''=u_0v''\mathcal A$ with relative error
$O(t+|u_0|+\delta_\epsilon t)$.
All zeros of $M'$ are simple and cross with sign $u_0$, so there
is at most one. This controls the auxiliary numerator. Actual returns
remain subject to the strict deficit conditions.

Suppose $M$ has two zeros in the comparison interval. The auxiliary
interval contains both zeros and extends to $t=0$. If $u_0,z_0$
were both negative, its only possible critical point would be a maximum, whereas
$M(0+)=\epsilon\kappa u_0z_0>0$. Such a function has at most one
positive zero. Hence
\begin{equation}\label{rw:finite:three-away-positive}
 u_0>0,\quad z_0>0,\quad
 z_0\le Cu_0t_*^{r-1}\le Cu_0t_0^\delta
\end{equation}
for a critical point $t_*$ between the two zeros. In particular,
after reducing $t_0$, $\xi_2=(u_0-z_0)/2>0$.

To recover $\tau_1$ with a uniform sign, work first at
$\epsilon=\nu=0$ on $z_0=0$. Let $q=q_l=Q_s$ be the common
incoming root. Its lower and upper equations are
\[
 \beta q^2+\bar a q=1,\qquad
 \bar m q^2-\bar c q=1.
\]
The outgoing intersections are $p_l=1/(\beta q)$ and
$P_s=1/(\bar m q)$. Therefore exactly
\begin{equation}\label{rw:finite:three-away-transverse}
 \tau_1=(\bar m-\beta)q,\qquad
 \delta_+=p_l-P_s=\frac{\bar m-\beta}{\beta\bar m q}.
\end{equation}
Since $u_0=\chi_+(\delta_+)$ with $\chi_+'>c$, these formulas
give $\tau_1=A_0u_0$ on this surface, where $A_0>c$ has bounded
jets. In particular the derivative in $\bar m-\beta$ is a
positive unit at the center, even when the actual ratio is used
as the remaining coordinate.

In the common inverse coordinates define, at positive scale,
$A_\epsilon=\int_0^1\partial_{u_0}
\tau_1(su_0,0,r,0;\epsilon)\,\dd s$.
The exact triangle identity makes the omitted constant zero.
Convergence of the lower strong intersections and the upper
low-layer jets to their zero-scale values gives
$|A_\epsilon-A_0|\le C\epsilon(1+|\log\epsilon|)$.
Single division in $z_0,\nu$ now gives
\[
 \tau_1=A_\epsilon u_0+B_\epsilon z_0+C_\epsilon\nu,\qquad
 A_\epsilon>c,\quad |B_\epsilon|+|C_\epsilon|\le C.
\]
Use \eqref{rw:finite:three-away-positive} and
$|\nu|\le\delta_\epsilon(u_0+z_0)$.
Choose $t_0$ with $C t_0^\delta$ small first, then the scale with
$C\delta_\epsilon$ small. The two errors are absorbed by
$A_\epsilon u_0/2$, and hence $\tau_1>0$. The lower and upper
intersection equations, together with the critical-point balance,
therefore determine the coefficient direction.
The same argument applied to $\bar m-\beta$ gives
$\bar m-\beta>0$: on the zero-scale surface in
\eqref{rw:finite:three-away-transverse} its quotient by $u_0$
is also a positive bounded unit, and it vanishes on the exact
positive-scale triangle axis. On the compact center part with
$r\ge1+\delta$, one has $\bar c=-\bar a<-c_\delta<0$;
the actual ratio coordinates preserve this margin on a smaller
neighborhood. Therefore the original primary generator has
$\sigma_1=\bar c(\bar m-\beta)<0$ on this nonresonant two-zero
branch. All $\nu$ terms in these divisions are $O(|\nu|)$ and
are absorbed by the stated forcing; no extra $\epsilon$ factor
is assigned to the $\nu$ dependence of the ratio.

The local calculation gives an upper bound of four closing zeros. It also identifies the
only branch on which the numerator can have two zeros. The signs on
that branch will be used at the already fixed regular seam in
Section~12; they do not yet determine the value of $\mathcal K$ at
that seam.

\begin{proposition}[Local count and the two-zero terminal branch]
\label{rw:finite:three-count}
On common small neighborhoods of every compact triangle-center
stratum, the closing equation has at most four isolated zeros on
\eqref{rw:finite:three-domain}, with multiplicity.
More precisely, on a noncenter fibre satisfying the local forcing
bound supplied by one closing zero,
\[
 Z(M)\le2,\qquad Z(\mathcal K)\le3,\qquad Z(F)\le4.
\]
If a terminal subinterval contains two zeros of $M$, counted with
multiplicity, orient the whole field by reflection if necessary so
that $r>1$. Then
\begin{equation}\label{rw:finite:three-terminal-signs}
 u_0>0,\qquad z_0>0,\qquad \xi_2>0,\qquad \tau_1>0,
 \qquad M(0+)>0,\qquad
 \mathcal K(t)\longrightarrow-\infty\quad(t\downarrow0).
\end{equation}
Near ratio one the same branch also has $T>0$, $W>0$ and
$\sigma_s<0$; away from one it has $\bar m>\beta$ and
$\sigma_1<0$. In particular $t=0$ is the end of the deficit
comparison interval in this branch. It need not be an end of the
positive-central return domain. If a seam $t_s$ in the comparison interval
is fixed away from the zeros and its tail contains exactly these two
zeros, then $M(t_s)>0$, $\mathcal K'_t(t_s)>0$, and
$\mathcal K'_p(t_s)<0$ in the outward coordinate
$p=P_s-\psi_+(t)$, $\psi_+'>0$.
These conclusions do not assume or provide a sign for
$\mathcal K(t_s)$.
\end{proposition}

\begin{proof}
Near one, \eqref{rw:finite:three-remainder} and
\eqref{rw:finite:three-column-norm} satisfy
Lemma~\ref{rw:finite:compensator-calculus}. Away from one, the
single crossing established above gives $Z(M)\le2$.
On the comparison interval, $M=tuz\,\mathcal K'$ with $tuz>0$,
and $F'=D_+'(u)G_+'(t)(1-e^{\mathcal K/\epsilon})$.
Two applications of Rolle's theorem with multiplicity therefore give
the displayed bounds. If there is no closing zero, the bound for $F$
is immediate. If forcing gives the exact center, there is no isolated
closing zero.

The away-from-one terminal signs were proved in
\eqref{rw:finite:three-away-positive} and
\eqref{rw:finite:three-away-transverse}. For the near-one signs,
first restrict to a smaller ratio interval so that its reflected
ratio remains in $|r-1|\le1/64$.
Fix a large $K_0$ and separate the cone
$|\xi_2|\le K_0|\kappa\xi_1|$.
On this cone, the exact positive column
\eqref{rw:finite:three-W-factor}, the small $\xi_2$ column,
and relative forcing give
\[
 \frac{W}{\kappa\xi_1}>c,\qquad
 |T/\xi_1|\le C|\kappa|,\qquad
 |\mathcal R_3'/\xi_1|\le C|\kappa|t^{1/2}.
\]
Thus the one-zero part of
Lemma~\ref{rw:finite:compensator-calculus} applies to $M/\xi_1$.
For $\kappa=0$ this cone is the reversible center. Neither case
permits the two-zero branch.

In the complementary cone choose $K_0$ large and then reduce the
coefficient box so that $T$ has the sign of $\xi_2$.
Two zeros of $M$ give a zero of $M'$ in the deficit interval, so
Lemma~\ref{rw:finite:three-critical-gaps} implies $u_0z_0>0$.
The case $\kappa=0$ would have $C_0=0$, contrary to
\eqref{rw:finite:compensator-signs}. For $\kappa>0$, $C_0>0$,
so that same lemma gives $W>0$ and $T>0$, hence $\xi_2>0$.
Equation~\eqref{rw:finite:three-directed-generators}, this cone and
forcing imply $\tau_1>0$ and $\sigma_s<0$.
The exact formula
\[
 W=\kappa u_0Q_*+2\epsilon u_0z_0\sigma_s r
\]
then excludes $u_0,z_0<0$: both terms would be negative.
Thus both deficits are positive.

Reflection exchanges the lower legs and replaces the actual upper
passage by its inverse. The new ratio is $1/r$, the new critical
function is $-\mathcal K$ composed with the inverse input, and its
derivative numerator is a nonzero signed unit times the old one.
It preserves zeros with multiplicity and the comparison interval.
Apply the preceding argument in that complete reflected coordinate
choice when $r<1$; no inequality is transported while keeping the
old entry coordinate fixed.

Positive endpoint deficits now place all sufficiently small $t>0$
in the comparison interval. Both complete central outputs must still
be positive for an actual return. At fixed coefficients, the lower
derivatives at $u_0,z_0$ are finite and nonzero, while
$v'\sim rt^{r-1}$.
This proves the endpoint limit for $\mathcal K$ in
\eqref{rw:finite:three-terminal-signs}.
Finally $M(0+)>0$ and exactly two zeros before the fixed seam,
counted with multiplicity, imply $M(t_s)>0$ by parity.
The positive multiplier and $dp/dt<0$ give the two derivative signs.
The remaining step is to determine $\mathcal K(t_s)$ from the
regular overlap. Proposition~\ref{rw:anchors:three-saddle} combines
that value with the derivative signs and endpoint parity to count
zeros on the whole union. The local estimate alone leaves the
four-zero closing bound.
\end{proof}

